% 第 IV 讲 9.1--9.4；原 PDF 260--283 页。
\SourcePage{260}
\SGAsection{9}{子拓扑斯与拓扑斯的粘合}
\SGAsubsection{9.1}{子拓扑斯与拓扑斯嵌入}
\begin{SGAresult}{定义}{9.1.1}
a) 若\Term{topos-morphism}{拓扑斯态射} $f:F\to E$ 的\Term{direct-image-functor}{正像函子} $f_*:F\to E$ 是\Term{fully-faithful-functor}{全忠实的}（即\Term{adjunction-morphism}{伴随态射} $f^*f_*\to\operatorname{id}_F$ 是\Term{isomorphism}{同构}），则称 $f$ 为\Term{topos-embedding}{嵌入}。

b) \Term{topos}{拓扑斯} $E$ 的一个\Term{subtopos}{子拓扑斯}，是 $E$ 的一个\Term{strictly-full-subcategory}{严格满子范畴} $E'$，它本身是\Term{topos}{拓扑斯}，且\Term{inclusion-functor}{包含函子} $\alpha:E'\to E$ 具有 $i_*$ 的形式，其中 $i:E'\to E$ 是\Term{topos-morphism}{拓扑斯态射}（即依 3.1.1，$\alpha$ 有一个\Term{left-exact-functor}{左正合的}\Term{left-adjoint}{左伴随} $i^*$）。此时称 $i$（它在唯一\Term{isomorphism}{同构}意义下确定）为\Term{subtopos}{子拓扑斯} $E'$ 到 $E$ 的\Term{inclusion-morphism}{包含态射}。
\end{SGAresult}
\begin{SGAresult}{注记}{9.1.2}
a) 显然，\Term{subtopos}{子拓扑斯}的\Term{inclusion-morphism}{包含态射}是\Term{topos-embedding}{嵌入}；而 $E$ 的一个\Term{subcategory}{子范畴} $E'$ 是\Term{subtopos}{子拓扑斯}，当且仅当它是某个\Term{topos-morphism}{拓扑斯态射} $g:F\to E$ 所\Term{correspondence}{对应}的\Term{direct-image-functor}{正像函子} $f_*$ 的\Term{essential-image}{本质像}，且 $g$ 是\Term{topos-embedding}{嵌入}。

b) 从定义立刻得到：\Term{topos-morphism}{拓扑斯态射} $f:F\to E$ 是\Term{topos-embedding}{嵌入}，当且仅当它\Term{isomorphism}{同构}于复合\Term{morphism}{态射}
\[
F\xrightarrow{g}E'\xrightarrow{i}E,
\]
其中 $g$ 是\Term{equivalence-of-topoi}{拓扑斯等价}，$i$ 是 $E$ 的一个\Term{subtopos}{子拓扑斯} $E'$ 的\Term{inclusion-morphism}{包含态射}。后者由上述条件唯一确定，正是 $f_*$ 的\Term{essential-image}{本质像}；前述\Term{resolution}{分解}也在唯一\Term{isomorphism}{同构}意义下唯一。当然，实际应用中宜借 $g$ 将 $F$ 与 $E$ 的\Term{subtopos}{子拓扑斯} $E'$ 等同（比较 3.4.1）。

\SourcePage{261}
c) 一个\Term{equivalence-of-topoi}{拓扑斯等价}（3.4）$u:E\xrightarrow{\sim}E_1$ 显然给出 $E$ 的\Term{subtopos}{子拓扑斯}\Term{set}{集合}与 $E_1$ 的\Term{subtopos}{子拓扑斯}\Term{set}{集合}之间的\Term{bijection}{双射}。事实上，$u_*$ 是\Term{equivalence-of-categories}{范畴等价}，故在 $E$ 的\Term{strictly-full-subcategory}{严格满子范畴}\Term{set}{集合}与 $E_1$ 的\Term{strictly-full-subcategory}{严格满子范畴}\Term{set}{集合}之间建立\Term{bijection}{双射}（$E'\subset E$ \Term{correspondence}{对应}于 $u_*|E'$ 的\Term{essential-image}{本质像} $E'_1$）；立刻看出，$E'$ 是 $E$ 的\Term{subtopos}{子拓扑斯}，当且仅当 $E'_1$ 是 $E_1$ 的\Term{subtopos}{子拓扑斯}。因此，研究\Term{subtopos}{子拓扑斯}时，可以归结到 $E=\widetilde C$ 的情形，其中 $C$ 是一个小\Term{site}{景}。在此情形，由第 II 讲 5.5，$E$ 的\Term{strictly-full-subcategory}{严格满子范畴} $E'$ 是\Term{subtopos}{子拓扑斯}，当且仅当\Term{inclusion-functor}{包含函子} $\alpha:E'\to E$ 有一个\Term{left-exact-functor}{左正合的}\Term{left-adjoint}{左伴随}（这本身已蕴含 $E'$ 是\Term{topos}{拓扑斯}）；这些\Term{subtopos}{子拓扑斯}的\Term{set}{集合}与 $C$ 上比给定\Term{topology}{拓扑} $T$ 更细的\Term{topology}{拓扑} $T'$ 的\Term{set}{集合}一一\Term{correspondence}{对应}，对每个这样的 $T'$，配以由 $T'$ 的\Term{sheaf}{层}组成的 $C^\sim$ 的\Term{strictly-full-subcategory}{严格满子范畴} $(C,T')^\sim$。这同时说明，对任意 $\mathcal U$-\Term{topos}{拓扑斯} $E$：$E$ 的\Term{strictly-full-subcategory}{严格满子范畴} $E'$ 是\Term{subtopos}{子拓扑斯}，当且仅当\Term{inclusion-functor}{包含函子} $\alpha:E'\to E$ 有一个\Term{left-exact-functor}{左正合的}\Term{left-adjoint}{左伴随}；$E$ 的\Term{subtopos}{子拓扑斯}所成的\Term{ordered-set}{有序集合} $\mathscr E(E)$ 是 $\mathcal U$-小的，且 $\mathscr E(E)$ 的每个\Term{subset}{子集}都有\Term{supremum}{上确界}和\Term{infimum}{下确界}。

d) 若 $E'$ 和 $E''$ 是\Term{topos}{拓扑斯} $E$ 的两个\Term{subtopos}{子拓扑斯}，则它们的交是 $E$ 的\Term{subtopos}{子拓扑斯}。事实上，因这个交是\Term{strictly-full-subcategory}{严格满的}，可归结到 $E=C^\sim$ 的情形。采用 b) 的记号，$E'$ 和 $E''$ \Term{correspondence}{对应}于 $C$ 上比 $T$ 更细的两个\Term{topology}{拓扑} $T',T''$，只须证明 $E'\cap E''$ 是\Term{topology}{拓扑} $T'''$ 的\Term{sheaf}{层}\Term{category}{范畴}，其中 $T'''$ 是 $T'$ 和 $T''$ 的\Term{supremum}{上确界}。下面给出证明（这个事实本应作为推论列在第 II 讲 4.4 或第 II 讲 5.5 之后！）。设 $E'''$ 是 $T'''$ 的\Term{sheaf}{层}\Term{category}{范畴}；它显然包含于 $E'$ 和 $E''$，故包含于 $E'\cap E''$。对 $E$ 的任意\Term{full-subcategory}{满子范畴} $D$，记 $t(D)$ 为使 $D$ 的元素都是\Term{sheaf}{层}的\Term{topology}{拓扑}中最细的一个（第 II 讲 2.2）。包含关系 $E'''\subset E'\cap E''\subset E',E''$ 显然蕴含
\[
t(E'),t(E'')\le t(E'\cap E'')\le t(E''').
\]
另一方面，由第 II 讲 4.4 I，有 $t(E')=T'$、$t(E'')=T''$、$t(E''')=T'''$，故以上不等式与 $T'''$ 作为\Term{supremum}{上确界}的定义蕴含 $T'''=t(E'\cap E'')$，从而 $E'\cap E''\subset E'''$，证毕。更一般地，这个论证说明：若\Term{topology}{拓扑} $T'$ 是任意一族\Term{topology}{拓扑} $(T_i)$ 的\Term{supremum}{上确界}，则 $T'$ 的\Term{sheaf}{层}\Term{category}{范畴} $E'$ 是各 $T_i$ 的\Term{sheaf}{层}\Term{category}{范畴} $E_i$ 的交。特别得到：
\end{SGAresult}
\begin{SGAresult}{命题}{9.1.3}
设 $E$ 是一个\Term{topos}{拓扑斯}。则对 $E$ 的任意一族\Term{subtopos}{子拓扑斯} $(E_i)_{i\in I}$，其交 $\bigcup_{i\in I}E_i$ 是 $E$ 的一个\Term{subtopos}{子拓扑斯}。
\end{SGAresult}
\begin{SGAresult}{命题}{9.1.4}
设 $i:E'\to E$ 是\Term{topos-embedding}{拓扑斯嵌入}，$F$ 是一个\Term{topos}{拓扑斯}。\Term{functor}{函子}
\[
f\longmapsto i\circ f:\operatorname{Homtop}(F,E')\longrightarrow\operatorname{Homtop}(F,E)\tag{9.1.4.1}
\]
\SourcePage{262}
是\Term{fully-faithful-functor}{全忠实的}，其\Term{essential-image}{本质像}由这样的\Term{topos-morphism}{拓扑斯态射} $g:F\to E$ 组成：$g_*$ 的\Term{essential-image}{本质像}包含于 $i_*$ 的\Term{essential-image}{本质像}。
\end{SGAresult}
\begin{proof}
考虑下列显然交换的图：
\[
\begin{tikzcd}
\operatorname{Homtop}(F,E')\arrow[r]\arrow[d,hook]&\operatorname{Homtop}(F,E)\arrow[d,hook]\\
\operatorname{Hom}(F,E')\arrow[r]&\operatorname{Hom}(F,E),
\end{tikzcd}
\]
其中竖\Term{arrow}{箭头}是\Term{functor}{函子} $h\mapsto h_*$，按 $\operatorname{Homtop}$ 的定义，它们是\Term{fully-faithful-functor}{全忠实的}（第 IV 讲 3.2）。另一方面，第二个横\Term{arrow}{箭头} $G\mapsto i_*G$ 是\Term{fully-faithful-functor}{全忠实的}，因为 $i_*$ 是\Term{fully-faithful-functor}{全忠实的}，所以第一个横\Term{arrow}{箭头}也如此。只剩证明 (9.1.4.1) 的\Term{essential-image}{本质像}的刻画。由公式 $(if)_*=i_*f_*$，条件的必要性显然；故只须证明，若 $g_*$ 的\Term{essential-image}{本质像}包含于 $i_*$ 的\Term{essential-image}{本质像}，即若可写成 $g_*=f_*i_*$，其中 $f_*:F\to E'$ 是一个\Term{functor}{函子}，则 $g$ 属于 (9.1.4.1) 的\Term{essential-image}{本质像}。这\Term{equivalence}{等价}于说 $f_*$ 有一个\Term{right-exact-functor}{右正合的}\Term{left-adjoint}{左伴随}。但显然 $f_*$ 有\Term{left-adjoint}{左伴随} $g^*i_*$，而这个\Term{functor}{函子}是\Term{left-exact-functor}{左正合的}，因为它是两个\Term{left-exact-functor}{左正合函子} $g^*$ 和 $i_*$ 的复合。
\end{proof}
关于 9.1.4 的一个逆命题，见练习 9.1.6。

取 $F$ 为\Term{punctual-topos}{点拓扑斯}（第 IV 讲 2.2），由 9.1.5 特别得到：
\begin{SGAresult}{推论}{9.1.5}
若 $f:E'\to E$ 是\Term{topos-embedding}{拓扑斯嵌入态射}，则相应的\Term{functor}{函子}
\[
\operatorname{Point}(f):\mathscr{P}\mathrm{oint}(E')\longrightarrow\mathscr{P}\mathrm{oint}(E)\tag{9.1.5.1}
\]
是\Term{fully-faithful-functor}{全忠实的}。
\end{SGAresult}
\begin{SGAresult}{练习}{9.1.6}
（\Term{topos-embedding}{拓扑斯嵌入}与 $2$-\Term{monomorphism}{单态射}。）设 $i:E'\to E$ 是\Term{topos-morphism}{拓扑斯态射}。

a) 若 $F,G$ 是 $E'$ 上方的两个\Term{topos}{拓扑斯}，定义一个\Term{canonical}{典范}\Term{functor}{函子}（采用 5.14 a) 的记号）
\[
f\longmapsto f*i:\operatorname{Homtop}_{E'}(F,G)\longrightarrow\operatorname{Homtop}_{E}(F,G).\tag{9.1.6.1}
\]
b) 假设 $i$ 作为 $\mathcal V$-$\mathcal U$-\Term{topos}{拓扑斯}的 $2$-\Term{category}{范畴}（3.3.2）中的 $1$-\Term{arrow}{箭头}，是一个\Term{two-monomorphism}{$2$-单态射}，即对属于 $\mathcal V$ 的每个 $\mathcal U$-\Term{topos}{拓扑斯} $F$，\Term{functor}{函子} (9.1.4.1) 都是\Term{fully-faithful-functor}{全忠实的}。证明：对任意一对 $\mathcal U$-\Term{topos}{拓扑斯} $F,G$（不必属于 $\mathcal V$），\Term{functor}{函子} (9.1.6.1) 都是\Term{fully-faithful-functor}{全忠实的}。

c) 证明：$i$ 是\Term{topos-embedding}{嵌入}，当且仅当 $i$ 是\Term{two-monomorphism}{$2$-单态射}。（证明充分性时，对\Term{induced-topos}{诱导拓扑斯} $E'_{/X'}$、$E'_{/Y'}$ 应用 b)，并使用 5.14 c)。）

\SourcePage{263}
d) 设 $g:F\to E$ 是\Term{topos-morphism}{拓扑斯态射}，并假设\Term{two-fiber-product}{$2$-纤维积}（5.11）$F'=F\mathbin{\mathop{\times}\limits_E^{2}}E'$ 存在（P. Deligne 已证明这个条件总成立）。设 $j:F'\to F$ 是\Term{projection}{投影}\Term{morphism}{态射}。证明：若 $i$ 是\Term{topos-embedding}{嵌入}，则 $j$ 也是。（使用 c)，并注意\Term{two-monomorphism}{$2$-单态射}在 $2$-\Term{base-change}{基变换}下的\Term{stable}{稳定}性是形式的。）当 $E'$ 是 $E$ 的\Term{subtopos}{子拓扑斯}，且 $i:E'\to E$ 是\Term{canonical}{典范}\Term{inclusion-morphism}{包含态射}时，$F'$ 在 $j_*$ 下的\Term{essential-image}{本质像}（它是 $F$ 的一个\Term{subtopos}{子拓扑斯}，不依赖于 $2$-\Term{fiber-product}{纤维积}构造中的不确定性）称为 $E$ 的\Term{subtopos}{子拓扑斯} $E'$ 在 $g:F\to E$ 下的\Term{inverse-image}{逆像}\Term{subtopos}{子拓扑斯}。

e) 采用 d) 末尾的记号，取 $F$ 中的\Term{subtopos}{子拓扑斯}——$E'$ 的\Term{inverse-image}{逆像}作为 $2$-\Term{fiber-product}{纤维积}。设 $X$ 是 $F$ 的\Term{object}{对象}。证明：$X$ 要属于 $F'$，必须满足下列条件：对 $F$ 的每个\Term{object}{对象} $U$，以 $g_U:F_{/U}\to E$ 表示由 $f$ 诱导的\Term{topos-morphism}{拓扑斯态射}，都有 $g_{U*}(X|U)\in\operatorname{Ob}E'$，其中 $X|U=(X\times U\xrightarrow{\operatorname{pr}_2}U)$ 是“$X$ 在 $U$ 上的\Term{restriction}{限制}”。（注意：若 $i_U:F_{/U}\to F$ 是\Term{canonical}{典范}包含，则 $g_U=g\circ i_U$，从而
\[
g_{U*}(X|U)=g_*(i_{U*}(X|U))=g_*(\mathscr{H}\mathrm{om}(U,X)),
\]
其中 $\mathscr{H}\mathrm{om}(U,X)$ 是下文 10.1 定义的\Term{object}{对象}。）

\textbf{问题。}反过来，$X$ 满足上述条件，是否就足以保证 $X$ 属于 $F'$？这\Term{equivalence}{等价}于问：由这些\Term{object}{对象} $X$ 组成的 $F$ 的\Term{strictly-full-subcategory}{严格满子范畴}是否是 $F$ 的一个\Term{subtopos}{子拓扑斯}？

f) 设 $X$ 是 $E$ 的\Term{object}{对象}，$X'=i^*(X)$，$i_{/X}:E'_{/X'}\to E_{/X}$ 是由 $i$ 诱导的\Term{topos-morphism}{拓扑斯态射}（5.10.1）。证明：若 $i$ 是\Term{topos-embedding}{嵌入}，则 $i_{/X}$ 也是。（使用 d) 和 5.11。）反过来，若 $(X_\alpha)_\alpha$ 是\Term{covering}{覆盖}\Term{final-object}{终对象}的一族\Term{object}{对象}，且对每个 $\alpha$，$i_{/X_\alpha}$ 都是\Term{topos-embedding}{嵌入}，则 $i$ 也是。
\end{SGAresult}
\begin{SGAresult}{练习}{9.1.7}
（\Term{subtopos}{子拓扑斯}与\Term{arrow}{箭头}的\Term{multiplicative-set}{乘法集}。）

a) 设 $i:E'\to E$ 是\Term{topos-morphism}{拓扑斯态射}，$S$ 是 $E$ 中满足 $i^*(u)$ 为\Term{isomorphism}{同构}的\Term{arrow}{箭头} $u$ 的\Term{set}{集合}；考虑由 $i^*$ 诱导的\Term{canonical}{典范}\Term{functor}{函子}（参见 [2，第 I 章] 或第 VI 讲）：
\[
\rho:E[S^{-1}]\longrightarrow E'.\tag{9.1.7.1}
\]
证明：$i$ 是\Term{topos-embedding}{嵌入}，当且仅当前述\Term{functor}{函子} $\rho$ 是\Term{equivalence-of-categories}{范畴等价}；此时 $S$ 既容许\Term{left-calculus-of-fractions}{左分式演算}，也容许\Term{right-calculus-of-fractions}{右分式演算}。（使用 [2，第 I 章 1.3]。）在此情形下，$i_*$ 的\Term{essential-image}{本质像}可以由 $S$ 恢复：它由所有满足如下条件的 $X\in\operatorname{Ob}E$ 组成：对 $S$ 中每个 $u:Y\to Y'$，
\[
\operatorname{Hom}(u,X):\operatorname{Hom}(Y',X)\longrightarrow\operatorname{Hom}(Y,X)
\]
都是\Term{bijection}{双射}。

b) 以下假设 $i$ 是\Term{topos-embedding}{嵌入}。证明：对任意\Term{topos}{拓扑斯} $F$，\Term{functor}{函子}
\[
f^*\longmapsto f^*i^*:\operatorname{Hom}(E',F)\longrightarrow\operatorname{Hom}(E,F)
\]
\SourcePage{264}
在下列两个\Term{category}{范畴}之间诱导\Term{equivalence-of-categories}{等价}：来自\Term{topos-morphism}{拓扑斯态射} $f:F\to E'$ 的 $f^*$（即\Term{left-exact-functor}{左正合}且有\Term{right-adjoint}{右伴随}者）的\Term{category}{范畴}；以及来自\Term{topos-morphism}{拓扑斯态射} $g:F\to E$、并把 $S$ 的\Term{object}{对象}变为\Term{isomorphism}{同构}的\Term{functor}{函子} $g^*:E\to F$ 的\Term{category}{范畴}。

c) 设 $S\subset\operatorname{Fl}E$，其中 $E$ 是一个 $\mathcal U$-\Term{topos}{拓扑斯}。证明：$S$ 可由 a) 的过程关联于 $E$ 的一个\Term{subtopos}{子拓扑斯} $E'$，当且仅当 $S$ 满足下列条件：
\begin{enumerate}[label={ST \arabic*)}]
\item $S$ 包含所有\Term{invertible}{可逆}\Term{arrow}{箭头}，且对复合与\Term{base-change}{基变换}\Term{stable}{稳定}。
\item 若复合 $vu$ 属于 $S$，则 $u\in S$ 当且仅当 $v\in S$。
\item 若 $u:X\to Y$ 满足：存在一个\Term{covering-family}{覆盖族} $Y_i\to Y$，$i\in I$，使得对所有 $i\in I$，$u_i:X_i=X\times_Y Y_i\to Y_i$ 都属于 $S$，则 $u\in S$。
\end{enumerate}
证明：这样得到 $E$ 的\Term{subtopos}{子拓扑斯} $E'$ 的\Term{set}{集合}与 $\operatorname{Fl}E$ 中满足 ST 1)、ST 2)、ST 3) 的\Term{subset}{子集} $S$ 的\Term{set}{集合}之间的一一\Term{correspondence}{对应}。（给定 $S$，给它配以比 $E$ 的\Term{canonical-topology}{典范拓扑} $T$ 更细的 $\mathcal U$-\Term{topology}{拓扑} $T'$；其\Term{covering-family}{覆盖族} $(X_i\to X)_{i\in I}$ 是那些满足如下条件的族：各 $X_i$ 的像的包含 $X'\hookrightarrow X$ 属于 $S$。证明 $T'$ 的\Term{sheaf}{层}\Term{category}{范畴}\Term{equivalence}{等价}于 $E$ 的某个\Term{subtopos}{子拓扑斯} $E'$，并使用第 II 讲 4.4 证明它给回 $S$。）

d) 证明：只要知道 $S$ 中由\Term{monomorphism}{单态射}组成的\Term{subset}{子集} $S_0$，便知道 $S$。证明：映射 $S\mapsto S_0$ 在 $\operatorname{Fl}E$ 中满足 c) 的 ST 1) 至 ST 3) 的\Term{subset}{子集} $S$ 的\Term{set}{集合}，与 $\operatorname{Fl}E$ 中由\Term{monomorphism}{单态射}组成且满足同样条件 ST 1) 至 ST 3) 的\Term{subset}{子集}的\Term{set}{集合}之间建立\Term{bijection}{双射}。

e) 设 $(E'_i)_{i\in I}$ 是 $E$ 的一族\Term{subtopos}{子拓扑斯}，$(S_i)_{i\in I}$ 是 $\operatorname{Fl}E$ 中相应的一族\Term{subset}{子集}。证明：$\bigcap_{i\in I}S_i$ 满足 c) 的 ST 1) 至 ST 3)，故\Term{correspondence}{对应}于 $E$ 的一个\Term{subtopos}{子拓扑斯} $E'$；后者是由各 $E_i$ 生成的\Term{subtopos}{子拓扑斯}，即各\Term{subtopos}{子拓扑斯} $E'_i$ 的\Term{supremum}{上确界}（参见 9.1.2 c)）。

f) 设 $A\subset\operatorname{Fl}E$。证明：在包含 $A$ 且满足 ST 1) 至 ST 3) 的 $\operatorname{Fl}E$ 的\Term{subset}{子集}中，存在最小的一个 $S$；$S$ \Term{correspondence}{对应}于最大的\Term{subtopos}{子拓扑斯} $E'$，使其\Term{correspondence}{对应}的 $\operatorname{Fl}E$ 的\Term{multiplicative-set}{乘法子集}包含 $S$。证明：$S$ 等于按如下方式递归构造的 $\operatorname{Fl}E$ 的\Term{subset}{子集} $A_n$（$n\ge0$）的并：$A_0=A$，$A_{n+1}$ 由下列三种类型 (i) 至 (iv) 之一的\Term{arrow}{箭头} $u:X\to Y$ 组成：
\begin{enumerate}[label=(\roman*)]
\item $u$ \Term{invertible}{可逆}，或是 $A_n$ 中两个\Term{arrow}{箭头}的复合，或由 $A_n$ 中某个\Term{arrow}{箭头}经\Term{base-change}{基变换}得到。
\item $u$ 是某个复合\Term{arrow}{箭头}的一个因子，且另一因子及该复合都属于 $A_n$。
\item $u$ 是 $A_n$ 中某个小\Term{arrow}{箭头}族的\Term{coproduct}{和}。
\SourcePage{265}
\item 存在\Term{epimorphism}{满态射} $Y'\to Y$，使得由 $u$ 经\Term{base-change}{基变换}得到的 $u'=X'=X\times_Y Y'\to Y'$ 属于 $A_n$。
\end{enumerate}
g) 采用 e) 的记号，设 $A$ 是各 $S_i$ 的并，$S$ 是由 $A$ 经 f) 的过程得到的 $\operatorname{Fl}E$ 的\Term{subset}{子集}，$E'$ 是 $E$ 的相应\Term{subtopos}{子拓扑斯}。证明：$E'$ 是各\Term{subtopos}{子拓扑斯} $E'_i$ 的交（即\Term{infimum}{下确界}）。
\end{SGAresult}
\begin{SGAresult}{练习}{9.1.7.2}
（\Term{topos-morphism}{拓扑斯态射}的\Term{canonical}{典范}\Term{resolution}{分解}。）

a) 设 $f:F\to E$ 是\Term{topos-morphism}{拓扑斯态射}。证明：在\Term{isomorphism}{同构}意义下可找到 $f$ 的\Term{resolution}{分解}
\[
F\xrightarrow{f'}E'\xrightarrow{i}E,\tag{9.1.7.3}
\]
其中 $f'$ 是\Term{conservative-morphism}{保守的}（即依 6.4.0，$f'^*$ 是\Term{conservative-functor}{保守的}，也即\Term{faithful-functor}{忠实的}），$i$ 是\Term{topos-embedding}{嵌入}，且这个\Term{resolution}{分解}在\Term{equivalence}{等价}意义下唯一（参见 9.1.4）。为此，引入 $E$ 中使 $f^*(u)$ 为\Term{isomorphism}{同构}的\Term{arrow}{箭头} $u$ 所成的\Term{set}{集合} $S_f\subset\operatorname{Fl}E$，证明 $S_f$ 满足 9.1.7 c) 的条件，取 $E'$ 为相应的\Term{subtopos}{子拓扑斯}，并由 (9.1.7.1) 定义 $f'^*$；依 9.1.7 b)，它\Term{correspondence}{对应}于一个\Term{topos-morphism}{拓扑斯态射} $f'$。

b) 由\Term{topos-embedding}{嵌入} $i$ 定义的 $E$ 的\Term{subtopos}{子拓扑斯}，称为\Term{topos-morphism}{拓扑斯态射} $f$ 的\Term{image-subtopos}{像子拓扑斯}。证明：它是在\Term{isomorphism}{同构}意义下使 $f$ 经由 $E_1$ \Term{resolution}{分解}的所有\Term{subtopos}{子拓扑斯} $E_1$ 中最小的一个，即依 9.1.4，它是包含 $f_*$ 的像（或 $f_*$ 的\Term{essential-image}{本质像}）的最小\Term{subtopos}{子拓扑斯}。$f$ 是\Term{topos-embedding}{嵌入}的充要条件，是 $f$ 诱导 $E$ 与它的像之间的一个\Term{equivalence}{等价} $f'$；$f$ 是\Term{conservative-morphism}{保守的}的充要条件，是它的像等于整个 $E$。

c) 假设 $f=\operatorname{Top}(f_0)$ 关联于\Term{sober-space}{清醒}\Term{topological-space}{拓扑空间}之间的\Term{continuous-map}{连续映射} $f_0:Y\to X$（4.1）。考虑 $f_0$ 的\Term{canonical}{典范}\Term{resolution}{分解}
\[
Y\xrightarrow{f'_0}X'=f_0(Y)\xrightarrow{i_0}X,
\]
其中 $i_0$ 是包含。证明：$f$ 的\Term{canonical}{典范}\Term{resolution}{分解}可与相应\Term{resolution}{分解}
\[
\operatorname{Top}(Y)\xrightarrow{\operatorname{Top}(f'_0)}\operatorname{Top}(X')\lhook\joinrel\xrightarrow{\operatorname{Top}(i_0)}\operatorname{Top}(X)
\]
等同。设 $Z$ 是 $X$ 中包含 $X'=f(Y)$ 的最小\Term{sober-space}{清醒}子\Term{space}{空间}，即由满足 $\overline{\{x\}}\cap X'$ 在 $\overline{\{x\}}$ 中\Term{dense-subset}{稠密}的点 $x\in X$ 组成的\Term{set}{集合}。（注意：若 $X$ 的点都是闭的，或若 $X$ \Term{locally-noetherian}{局部诺特}且 $X'$ 是 $X$ 的一个\Term{constructible-subset}{可构造子集}，则 $Z=X'$。）则 $f$ 是\Term{conservative-morphism}{保守的}，即 $F=\operatorname{Top}(Y)$ 在 $f=\operatorname{Top}(f_0)$ 下的\Term{image-subtopos}{像子拓扑斯}等于 $E=\operatorname{Top}(X)$ 本身，当且仅当 $Z=X$，即当且仅当 $f(Y)$ 在 $X$ 中\Term{very-dense-subset}{很稠密}（EGA IV 10.13）。这说明，在何种程度上可以将\Term{topos}{拓扑斯}的\Term{conservative-morphism}{保守态射}概念，可以将\Term{topos}{拓扑斯}的\Term{conservative-morphism}{保守态射}概念视为\Term{topological-space}{拓扑空间}的\Term{surjectivity}{满射}\Term{continuous-map}{连续映射}概念的自然推广。
\SourcePage{266}
d) 设
\[
\begin{tikzcd}
Y\arrow[d]&Y'\arrow[l]\arrow[d]\\
X&X'\arrow[l]
\end{tikzcd}
\]
是\Term{topological-space}{拓扑空间}的\Term{cartesian-diagram}{笛卡尔图}。若 $Y$ 和 $X'$ 是\Term{discrete-space}{离散空间}，从而 $Y'$ 也是\Term{discrete-space}{离散的}，则相应的\Term{topos}{拓扑斯}图是 $2$-\Term{cartesian-diagram}{笛卡尔的}（5.11）。（注意：不再假设 $Y,X'$ \Term{discrete}{离散}时，即使 $Y,X'$ 是\Term{hausdorff-space}{分离空间} $X$ 的子\Term{space}{空间}，也可有反例，见 9.1.10 c)。）由此给出一个 $2$-\Term{cartesian-diagram}{笛卡尔}的\Term{topos}{拓扑斯}图
\[
\begin{tikzcd}
F\arrow[d,"f"']&F'\arrow[l]\arrow[d,"f'"]\\
E&E'\arrow[l]
\end{tikzcd}
\]
使得 $f$ 是\Term{conservative-morphism}{保守的}，但 $f'$ 不是\Term{conservative-morphism}{保守的}；更确切地，$E'$ 为\Term{punctual-topos}{点拓扑斯}，$F'$ 为\Term{empty-topos}{空拓扑斯}（2.2）。（取 $Y$ 使其在 $X$ 中的像很\Term{dense}{稠密}而不等于 $X$，再取 $X'$ 为不在该像中的一个点。）

e) 设 $(E_i)_{i\in I}$ 是\Term{topos}{拓扑斯} $E$ 的一族\Term{subtopos}{子拓扑斯}。证明：$E$ 是各 $E_i$ 的\Term{supremum}{上确界}，当且仅当\Term{topos-embedding}{嵌入态射}族 $E_i\to E$ 是\Term{conservative-family}{保守的}（6.4.0）。
\end{SGAresult}
\begin{SGAresult}{练习}{9.1.8}
（\Term{subtopos}{子拓扑斯}与 $E$ 的点；\Term{topological-space}{拓扑空间}的情形。）设 $E$ 是一个\Term{topos}{拓扑斯}。

a) 设 $Z$ 是 $\mathscr{P}\mathrm{oint}(E)$ 的一个\Term{full-subcategory}{满子范畴}（或\Term{equivalence}{等价}地，是 $\operatorname{Ob}\mathscr{P}\mathrm{oint}(E)$ 的一个\Term{subset}{子集}）。证明：被\Term{correspondence}{对应}于 $z\in\operatorname{Ob}Z$ 的\Term{fiber-functor}{纤维函子}变为\Term{bijection}{双射}的 $E$ 中的\Term{arrow}{箭头}所成的\Term{set}{集合} $S(Z)$ 满足 9.1.7 c) 的条件，因而定义 $E$ 的一个\Term{subtopos}{子拓扑斯} $E_Z$。证明 $E_Z$ 有\Term{enough-points}{足够多的点}，且 $Z$ \Term{equivalence}{等价}于 $\mathscr{P}\mathrm{oint}(E_Z)$ 的一个\Term{full-subcategory}{满子范畴}。证明 $E$ 中每个有\Term{enough-points}{足够多的点}的\Term{subtopos}{子拓扑斯} $E'$ 都可以由前述过程得到。（取 $Z$ 为\Term{functor}{函子} (9.1.5.1) 的\Term{essential-image}{本质像}。）证明：这样得到一个\Term{bijection}{双射}，一方是 $E$ 中有\Term{enough-points}{足够多的点}的\Term{subtopos}{子拓扑斯} $E'$ 的\Term{set}{集合}，另一方是 $\mathscr{P}\mathrm{oint}(E)$ 中对\Term{direct-factor}{直因子}（第 I 讲 10）及小\Term{filtered}{滤过}\Term{projective-limit}{投射极限}（参见 7.3）\Term{stable}{稳定}的\Term{strictly-full-subcategory}{严格满子范畴} $Z$ 的\Term{set}{集合}的某个\Term{subset}{子集}。（问题：这样得到的是 $\mathscr{P}\mathrm{oint}(E)$ 的哪些\Term{subcategory}{子范畴}？借助 9.1.2 c) 的\Term{correspondence}{对应}，这本质上是 6.5.3 中已经提出的问题的另一种形式。注意，除刚才指出的必要条件之外，还存在关于 $\operatorname{Point}(E)$ 的\Term{topology}{拓扑}性质、例如清醒性等的更细微条件，参见 e)。）

\SourcePage{267}
b) 设 $(Z_i)_{i\in I}$ 是 $\mathscr{P}\mathrm{oint}(E)$ 的一族\Term{full-subcategory}{满子范畴}，$Z$ 是各 $Z_i$ 的并所成的\Term{full-subcategory}{满子范畴}。证明：a) 中定义的相应\Term{subtopos}{子拓扑斯} $E_Z$ 是由各 $E_{Z_i}$ 生成的 $E$ 的\Term{subtopos}{子拓扑斯}。（使用 9.1.7 e)。）由此推出：由一族有\Term{enough-points}{足够多的点}的\Term{subtopos}{子拓扑斯}生成的\Term{subtopos}{子拓扑斯}也有\Term{enough-points}{足够多的点}。

c) 设 $f:E'\to E$ 是\Term{topos-embedding}{拓扑斯嵌入}。证明映射 $U\mapsto f^*(U)$
\[
\operatorname{Ouv}(E)\longrightarrow\operatorname{Ouv}(E')\tag{9.1.8.1}
\]
是\Term{surjectivity}{满射}，并由此推出：$f$ 在点的\Term{isomorphism}{同构}类上诱导的映射
\[
\operatorname{Point}(E')\longrightarrow\operatorname{Point}(E)\tag{9.1.8.2}
\]
对 7.8 a) 所定义的\Term{topology}{拓扑}而言，诱导 $\operatorname{Point}(E')$ 到 $\operatorname{Point}(E)$ 的一个子\Term{space}{空间}的\Term{homeomorphism}{同胚}。（\Term{injectivity}{单射}性使用 9.1.5。）

d) 设 $f:X'\to X$ 是\Term{continuous-map}{连续映射}，从而有\Term{topos-morphism}{拓扑斯态射}（4.1）
\[
\operatorname{Top}(f):\operatorname{Top}(X')\longrightarrow\operatorname{Top}(X).
\]
证明：后者是\Term{topos-embedding}{嵌入}，当且仅当 $X'$ 的\Term{topology}{拓扑}是 $X$ 的\Term{topology}{拓扑}在 $f$ 下的\Term{inverse-image}{逆像}（即当 $X'$ 是 Kolmogoroff \Term{space}{空间}，例如\Term{sober-space}{清醒空间}时，当且仅当 $f$ 诱导 $X'$ 到 $X$ 的一个子\Term{space}{空间}的\Term{homeomorphism}{同胚}）。（必要性使用 c)，充分性使用 9.1.7.2 c)。）

e) 设 $X$ 是\Term{sober-space}{清醒}\Term{topological-space}{拓扑空间}，$E=\operatorname{Top}(X)$。使用 a) 和 d)，建立 $X$ 的\Term{sober-space}{清醒}子\Term{space}{空间} $X'$ 的\Term{set}{集合}（即 $X$ \Term{separated}{分离}时，$X$ 的所有\Term{subset}{子集}的\Term{set}{集合}）与 $E$ 中有\Term{enough-points}{足够多的点}的\Term{subtopos}{子拓扑斯}的\Term{set}{集合}之间的\Term{order-isomorphism}{有序集合同构} $X'\mapsto E_{X'}$。证明 $E_{X'}$ \Term{equivalence}{等价}于 $\operatorname{Top}(X')$。若 $(X'_i)_{i\in I}$ 是 $X$ 的一族\Term{sober-space}{清醒}子\Term{space}{空间}，证明：存在包含各 $X'_i$ 的最小\Term{sober-space}{清醒}子\Term{space}{空间} $X'$；当 $I$ \Term{finite}{有限}或 $X$ \Term{separated}{分离}时，它等于各 $X'_i$ 的并；而 $E_{X'}$ 是由各 $E_{X'_i}$ 生成的 $E$ 的\Term{subtopos}{子拓扑斯}。（使用 b)。）

f) 给出一个形如 $\operatorname{Top}(X)$ 的\Term{topos}{拓扑斯} $E$ 及其\Term{subtopos}{子拓扑斯} $E'$，使 $E'$ 没有\Term{enough-points}{足够多的点}，甚至\Term{nonempty}{非空}（2.2）却没有点。（取 $E=\widehat U$、$E'=U^\sim$，其中 $U$ 是 7.4 中定义的\Term{ordered-set}{有序集合}。更好的是取 9.1.10 b) 的例子，它说明 $X$ 可取任意\Term{nonempty}{非空}且没有\Term{isolated-point}{孤立点}的\Term{compact-space}{紧空间}。）
\end{SGAresult}
\begin{SGAresult}{练习}{9.1.9}
（由\Term{ordered-set}{有序集合}定义的\Term{topos}{拓扑斯}的情形。）

a) 使用 7.8 b)、c) 的基本方法：它给出如下两方之间的\Term{correspondence}{对应}，一方是由其\Term{open-object}{开对象}生成的\Term{topos}{拓扑斯}及这些\Term{topos}{拓扑斯}的\Term{morphism}{态射}；另一方是具有任意\Term{supremum}{上确界}、\Term{finite}{有限}\Term{infimum}{下确界}，且二元\Term{infimum}{下确界}对任意\Term{supremum}{上确界}\Term{separable}{可分}配的\Term{ordered-set}{有序集合}，
\SourcePage{268}
其间的\Term{morphism}{态射}则是与\Term{finite}{有限}\Term{infimum}{下确界}及任意\Term{supremum}{上确界}交换的递增映射。

b) 设 $E$ 是\Term{topos}{拓扑斯}，$f:E'\to E$ 是\Term{topos-embedding}{拓扑斯嵌入}。证明：对 $E$ 中任意\Term{generating-family}{生成族} $(X_i)_{i\in I}$，族 $(f^*(X_i))_{i\in I}$ 也是\Term{generating-family}{生成族}。由此推出：若 $E$ 由其\Term{open-object}{开对象}生成，则 $E'$ 也如此。（使用 9.1.8 c)。）

c) 设 $f:E'\to E$ 是由各自\Term{open-object}{开对象}生成的\Term{topos}{拓扑斯}之间的\Term{morphism}{态射}，关联于相应\Term{ordered-set}{有序集合}上的\Term{morphism}{态射} $g:\mathscr O\to\mathscr O'$。证明：$f$ 是\Term{topos-embedding}{嵌入}，当且仅当 $g$ \Term{surjectivity}{满射}，且 $g$ 在\Term{order-relation}{序关系}的图上也\Term{surjectivity}{满射}，即 $\mathscr O'$ 的序是 $\mathscr O$ 的序的商。

d) 设 $E$ 是 a) 中形如 $\mathscr O^\sim$ 的\Term{topos}{拓扑斯}。由 b)、c) 推出：$E$ 的\Term{subtopos}{子拓扑斯}\Term{set}{集合}与 $\mathscr O$ 中具有下述性质的\Term{equivalence-relation}{等价关系} $R$ 的\Term{set}{集合}一一\Term{correspondence}{对应}；这些性质类似于 9.1.7 c) 的 ST 1) 至 ST 3)：
\begin{enumerate}[label={Q0 \arabic*)}]
\item 关系 $R$ 与\Term{finite}{有限}\Term{infimum}{下确界}相容，亦即：若 $U,U'$ 关于 $R$ \Term{equivalence}{等价}，则对每个 $V\in\mathscr O$，$U\cap V$ 与 $U'\cap V$ 也\Term{equivalence}{等价}。
\item 关系 $R$ 与任意\Term{supremum}{上确界}相容，亦即：若对所有 $i\in I$，$U_i$ 与 $U'_i$ \Term{equivalence}{等价}，则 $\sup_iU_i$ 与 $\sup_iU'_i$ \Term{equivalence}{等价}。
\end{enumerate}
证明：若 $E$ 的\Term{subtopos}{子拓扑斯} $E'_i$ \Term{correspondence}{对应}于\Term{equivalence-relation}{等价关系} $R_i$，则由各 $E'_i$ 生成的\Term{subtopos}{子拓扑斯}\Term{correspondence}{对应}于各 $R_i$ 的交关系。
\end{SGAresult}
\begin{SGAresult}{练习}{9.1.10}
（子\Term{space}{空间}的交；新的无\Term{punctual-topos}{点拓扑斯}；\Term{subtopos}{子拓扑斯}的补的不存在。）设 $X$ 是\Term{topological-space}{拓扑空间}，$E=\operatorname{Top}(X)=\operatorname{Ouv}(X)^\sim$（2.1）。

a) 在 $\operatorname{Ouv}(X)$ 上定义关系 $R$：$(U,U')\in R$ 当且仅当 $U\cap U'$ 在 $U$ 和 $U'$ 中都\Term{dense-subset}{稠密}。证明 $R$ 是满足 9.1.9 d) 的 Q0 1)、Q0 2) 的\Term{equivalence-relation}{等价关系}，故定义 $E$ 的一个\Term{subtopos}{子拓扑斯} $E'$。证明：对于点 $x\in X$，相应的点属于 $\mathscr{P}\mathrm{oint}(E')$ 的\Term{essential-image}{本质像}，当且仅当 $x$ 是\Term{thick-point}{厚点}（épais），即属于 $\overline{\{x\}}$ 的内部的\Term{closure}{闭包}。证明 $E'$ 是\Term{empty-topos}{空拓扑斯}（2.2），当且仅当 $X$ 为空。

b) 若 $X$ 清醒，证明\Term{category}{范畴} $\mathscr{P}\mathrm{oint}(E')$ \Term{equivalence}{等价}于由 $X$ 的\Term{thick-point}{厚点}构成的有序\Term{subset}{子集}（关于\Term{specialization-order}{特殊化关系}）所定义的\Term{category}{范畴}。证明：若 $X$ 的点都是闭的（括号中的相应情形：若 $X$ \Term{noetherian}{诺特}），则 $X$ 的\Term{thick-point}{厚点}就是 $X$ 的\Term{isolated-point}{孤立点}，即 $\{x\}$ 为\Term{open-subset}{开集}的点（相应地，就是 $X$ 的\Term{maximal-point}{极大点}）。由此推出：若 $X$ 是\Term{nonempty}{非空}、无\Term{isolated-point}{孤立点}的\Term{hausdorff-space}{分离}\Term{topological-space}{拓扑空间}，则 $E=\operatorname{Top}(X)$ 的\Term{subtopos}{子拓扑斯} $E'$ 是一个非“空”且没有点的\Term{topos}{拓扑斯}。

c) 设 $X'$ 是 $X$ 的\Term{subset}{子集}，$E_{X'}$ 是由 9.1.8 a) 的过程关联于 $X'$ 的 $E$ 的\Term{subtopos}{子拓扑斯}。证明：若 $X'$ \Term{dense}{稠密}，则 $E_{X'}$ 包含 $E'$。证明：若 $X$ 清醒，且 $X',X''$ 是 $X$ 的两个\Term{sober-space}{清醒}\Term{subset}{子集}，则 $X'\cap X''$ 清醒，且
\SourcePage{269}
$E_{X'\cap X''}\subset E_{X'}\cap E_{X''}$，而这个包含可以是严格的\footnote{但关于 $X$ \Term{locally-noetherian}{局部诺特}的情形，见下文 9.1.11 b)。}。（取一个\Term{nonempty}{非空}的\Term{hausdorff-space}{分离空间} $X$，它有两个不交的\Term{dense-subset}{稠密子集} $X',X''$。）证明前述包含为等号，当且仅当\Term{subtopos}{子拓扑斯} $E_{X'}\cap E_{X''}$ 有\Term{enough-points}{足够多的点}。证明若 $X'$ 或 $X''$ 是 $X$ 的\Term{locally-closed-subset}{局部闭子集}，则确实如此。

d) 设 $X'$ 是 $X$ 的一个\Term{subset}{子集}，$X''$ 是 $X'$ 的补集。$E$ 的\Term{subtopos}{子拓扑斯} $F$ 包含所有 $x\in X''$ \Term{correspondence}{对应}的 $E_{\{x\}}$，当且仅当它包含 $E_{X''}$（参见 9.1.8 b)）。由此推出：若 $X''$ 是 $X$ 的\Term{locally-closed-subset}{局部闭子集}的并（例如 $X''$ 的点都是闭的），则在满足 $F\cap E_{X'}$ 为\Term{empty-topos}{空子拓扑斯} $\operatorname{Top}(\varnothing)$ 的 $E$ 的\Term{subtopos}{子拓扑斯} $F$ 的\Term{set}{集合}中，只有当 $E_{X'}\cap E_{X''}$ 是\Term{empty-topos}{空子拓扑斯}时，才有最大元素（此时它值得称为 $E_{X'}$ 在 $E$ 中的\Term{weak-complement-subtopos}{弱补子拓扑斯}，参见 9.1.13）；当 $X\ne\varnothing$ 且 $X',X''$ 都在 $X$ 中\Term{dense}{稠密}，而 $X,X',X''$ 清醒时，这个条件不满足，参见 c)。证明：这个条件不满足时，在 $E$ 的\Term{subtopos}{子拓扑斯}所成的\Term{ordered-set}{有序集合}中，二元\Term{infimum}{下确界}对任意\Term{supremum}{上确界}不满足分配律。（执笔者不知道它对\Term{finite}{有限}\Term{supremum}{上确界}是否满足分配律；但参见 9.1.11 e) 和 9.1.12 d)。）
\end{SGAresult}
\begin{SGAresult}{练习}{9.1.11}
（\Term{locally-noetherian-topos}{局部诺特拓扑斯}的\Term{subtopos}{子拓扑斯}；\Term{noetherian-space}{诺特空间}的情形。）

a) 设 $f:E'\to E$ 是\Term{topos-embedding}{拓扑斯嵌入}。证明：若 $X$ 是 $E$ 的\Term{prenoetherian-object}{预诺特对象}（即依第 VI 讲 1，$x$ 的每个递增\Term{subobject}{子对象}序列最终\Term{stable}{稳定}），则 $X'=f^*(X)$ 是 $E'$ 的\Term{prenoetherian-object}{预诺特对象}。（引入\Term{induced-topos}{诱导拓扑斯} $E'_{/X'}$ 和 $E_{/X}$，并使用 9.1.6 e)，先归结到 $X$ 是 $E$ 的\Term{final-object}{终对象} $e$、从而 $X'$ 是 $E'$ 的\Term{final-object}{终对象} $e'$ 的情形。再注意 $e'$ 的每个递增\Term{subobject}{子对象}序列，经 $f_*$ 给出 $f_*(e')=e$ 的递增\Term{subobject}{子对象}序列。）由此推出：若 $E$ 有一个由\Term{prenoetherian-object}{预诺特对象}组成的\Term{generating-family}{生成族}（相应地，若 $E$ 是\Term{noetherian-topos}{诺特拓扑斯}，或\Term{locally-noetherian-topos}{局部诺特拓扑斯}，第 VI 讲 2），则 $E'$ 满足同样的条件。（使用 9.1.9 b)。）

b) 设 $E$ 是\Term{locally-noetherian-topos}{局部诺特拓扑斯}（第 VI 讲 2）。证明：$E$ 的每个\Term{subtopos}{子拓扑斯}都形如 $E_Z$（9.1.8 a)），其中 $Z$ 是 $\mathscr{P}\mathrm{oint}(E)$ 的适当\Term{full-subcategory}{满子范畴}。（使用 a)，以及 Deligne 的定理：每个\Term{locally-noetherian-topos}{局部诺特拓扑斯}都有\Term{enough-points}{足够多的点}。）由此推出：若 $E=\operatorname{Top}(X)$，其中 $X$ 是清醒的\Term{locally-noetherian-space}{局部诺特}\Term{topological-space}{拓扑空间}，则 $Z\mapsto E_Z$ 是 $X$ 的\Term{sober-space}{清醒}子\Term{space}{空间}所成\Term{ordered-set}{有序集合}到 $X$ 的\Term{subtopos}{子拓扑斯}所成\Term{ordered-set}{有序集合}的\Term{order-isomorphism}{同构}。（使用 9.1.8 e)。）\Term{sober-space}{清醒}子\Term{space}{空间} $Z_i$ 的交（必然清醒，且）定义各\Term{subtopos}{子拓扑斯} $E_{Z_i}$ 的交（与一般情形中可能发生的现象不同，9.1.10 c)）。

c) 设 $X$ 是清醒的\Term{locally-noetherian-space}{局部诺特空间}，$X'$ 是 $X$ 的\Term{subset}{子集}，$X''$ 为其补集。证明下列条件\Term{equivalence}{等价}：
\SourcePage{270}
\begin{enumerate}[label=(\roman*)]
\item $X'$ \Term{constructible-subset}{可构造}。
\item $X'$ 和 $X''$ 都清醒。
\item $X'$ 清醒，且 $E=\operatorname{Top}(X)$ 的\Term{subtopos}{子拓扑斯} $E_{X'}\simeq\operatorname{Top}(X')$ 有一个“\Term{complementary-subtopos}{补子拓扑斯}”（参见 9.1.13 e)）。
\item $E=\operatorname{Top}(X)$ 的\Term{subtopos}{子拓扑斯} $E_{X'}\simeq\operatorname{Top}(X')$ 与 $E_{X''}\simeq\operatorname{Top}(X'')$ 的交是 $E$ 的“\Term{empty-topos}{空}”\Term{subtopos}{子拓扑斯}。
\end{enumerate}
证明：若 $E',E''$ 是 $E$ 的两个\Term{subtopos}{子拓扑斯}，则它们互为\Term{weak-complement-subtopos}{弱补}，当且仅当它们互为\Term{complementary-subtopos}{补}（9.1.13）。

d) 设 $X$ 是清醒的\Term{noetherian-space}{诺特空间}。证明下列条件\Term{equivalence}{等价}：
\begin{enumerate}
\item[(i)] $X$ 的每个\Term{sober-space}{清醒}子\Term{space}{空间} $X'$ 都\Term{constructible-subset}{可构造}，即依 c)，其补集清醒。
\item[(i bis)] 对 $X$ 的每个\Term{maximal-point}{极大点} $x$，$X-\{x\}$ 清醒，即 $\{x\}$ \Term{constructible-subset}{可构造}。
\item[(ii)] $X$ \Term{finite}{有限}。
\item[(iii)] $\operatorname{Top}(X)$ 的每个\Term{subtopos}{子拓扑斯}都有\Term{complementary-subtopos}{补子拓扑斯}（参见 9.1.13）。
\item[(iv)] 在 $\operatorname{Top}(X)$ 的\Term{subtopos}{子拓扑斯}所成\Term{ordered-set}{有序集合}中，二元\Term{infimum}{下确界}对任意\Term{supremum}{上确界}满足分配律。
\end{enumerate}
（借助 b) 中证明的交的公式，使用 9.1.10 d) 的论证。）

e) 设 $X$ 是\Term{locally-noetherian-space}{局部诺特空间}。证明：在 $\operatorname{Top}(X)$ 的\Term{subtopos}{子拓扑斯}所成\Term{ordered-set}{有序集合}中，二元\Term{infimum}{下确界}对\Term{finite}{有限}\Term{supremum}{上确界}满足分配律。（归结到 $X$ 清醒的情形，并使用 b)。）

f) 对清醒的\Term{locally-noetherian-space}{局部诺特空间} $X$ 的每个\Term{sober-space}{清醒}子\Term{space}{空间} $Z$，以 $T_Z$ 表示 $X$ 的\Term{open-subset}{开集}\Term{category}{范畴} $\operatorname{Ouv}(X)$ 上的\Term{topology}{拓扑}，使一族包含 $U_i\to U$ 为\Term{covering-family}{覆盖族}，当且仅当 $U\cap Z=\bigcup_iU_i\cap Z$。证明：$Z\mapsto T_Z$ 在 $X$ 的\Term{sober-space}{清醒}子\Term{space}{空间} $Z$ 的\Term{set}{集合}，与 $\operatorname{Ouv}(X)$ 上比\Term{canonical-topology}{典范拓扑} $T_X$ 更细的\Term{topology}{拓扑}的\Term{set}{集合}之间建立\Term{bijection}{双射}（它反转包含关系）。（使用 b) 和 9.1.2 c)。）

g) 设 $X$ 是\Term{sober-space}{清醒}\Term{topological-space}{拓扑空间}。$X$ 的\Term{subset}{子集} $X'$ 清醒，当且仅当其补集是\Term{locally-closed-subset}{局部闭子集}的并；若 $X$ \Term{locally-noetherian}{局部诺特}，这也表示 $X'$ “\Term{proconstructible-subset}{pro-可构造}”（EGA IV 1.9.4）。证明：$X$ 上存在\Term{topology}{拓扑} $T_{\mathrm{cons}'}$，使得 $X_{\mathrm{cons}'}=(X,T_{\mathrm{cons}'})$ 的\Term{closed-subset}{闭子集}恰为 $X$ 的\Term{sober-space}{清醒}\Term{subset}{子集}。现在假设 $X$ \Term{locally-noetherian}{局部诺特}，则 $X_{\mathrm{cons}'}$ 正是 EGA IV 9.1.13 的\Term{space}{空间} $X_{\mathrm{cons}}$。由 b)、c) 推出：$\operatorname{Top}(X)$ 的\Term{subtopos}{子拓扑斯}所成\Term{ordered-set}{有序集合} $\Sigma$ \Term{isomorphism}{同构}于 $X_{\mathrm{cons}}$ 的\Term{closed-subset}{闭子集}所成\Term{ordered-set}{有序集合}，且在此\Term{correspondence}{对应}下，有\Term{complementary-subtopos}{补}的\Term{subtopos}{子拓扑斯}\Term{correspondence}{对应}于 $X_{\mathrm{cons}}$ 中既开又闭的\Term{subset}{子集}。证明 $X$ \Term{local-scheme}{局部}紧，且它紧当且仅当 $X$ \Term{noetherian}{诺特}，即\Term{topos}{拓扑斯} $\operatorname{Top}(E)$ \Term{noetherian}{诺特}。

\SourcePage{271}
h) 设 $E$ 是\Term{locally-noetherian-topos}{局部诺特拓扑斯}。假设它具有 9.1.14 b) 所述的性质（这个条件也许总满足……）。证明：当 $E'$ 遍历 $E$ 的\Term{subtopos}{子拓扑斯}\Term{set}{集合} $\Sigma$ 时，\Term{canonical}{典范}\Term{functor}{函子} $\mathscr{P}\mathrm{oint}(E')\to\mathscr{P}\mathrm{oint}(E)$ 的\Term{essential-image}{本质像}在 $E$ 的点的\Term{isomorphism}{同构}类\Term{set}{集合} $X=\operatorname{Point}(E)$ 中定义一族\Term{subset}{子集}，它恰是 $X$ 上某个\Term{topology}{拓扑} $T_{\mathrm{cons}}$ 的\Term{closed-subset}{闭子集}族；且由此得到\Term{ordered-set}{有序集合} $\Sigma$ 与 $X_{\mathrm{cons}}$ 的\Term{closed-subset}{闭子集}所成\Term{ordered-set}{有序集合}之间的\Term{order-isomorphism}{同构}。由此推出：在 $\Sigma$ 中，二元\Term{supremum}{上确界}对任意\Term{infimum}{下确界}满足分配律。证明：$X_{\mathrm{cons}}$ 本身虽不必是 $\mathcal U$-小的，但其关联\Term{sober-space}{清醒空间}（4.2.1）是 $\mathcal U$-小的。
\end{SGAresult}
\begin{SGAresult}{练习}{9.1.12}
（\Term{finite-topos}{有限拓扑斯}。）称\Term{topos}{拓扑斯} $E$ 为\Term{finite-topos}{有限的}，若它\Term{equivalence}{等价}于形如 $\widehat C$ 的\Term{topos}{拓扑斯}，其中 $C$ 是\Term{finite-category}{有限范畴}。

a) （\Term{correspondence}{对应}。）设 $\mathcal V$ 是满足 $\mathcal U\in\mathcal V$ 的\Term{universe}{宇宙}。以 $(\mathrm{Karfiness})$ 表示 $(\mathrm{Cat})_{\mathcal V}$ 的\Term{full-two-subcategory}{满 $2$-子范畴}，由属于 $\mathcal V$、为\Term{karoubian-category}{Karoubi 范畴}且\Term{equivalence}{等价}于\Term{finite-category}{有限范畴}的\Term{category}{范畴}组成；以 $(\mathrm{Topfin})$ 表示 $(\mathcal V\text{-}\mathcal U\text{-}\mathrm{Top})$（3.3.1）的\Term{full-two-subcategory}{满 $2$-子范畴}，由属于 $\mathcal V$ 的\Term{finite}{有限} $\mathcal U$-\Term{topos}{拓扑斯}组成。证明有互为\Term{quasi-inverse}{拟逆}的 $2$-\Term{equivalence}{等价}：
\[
C\longmapsto\widehat C:(\mathrm{Karfiness})\longrightarrow(\mathrm{Topfin}),\tag{9.1.12.1}
\]
\[
E\longmapsto\mathscr{P}\mathrm{oint}(E):(\mathrm{Topfin})\longrightarrow(\mathrm{Karfiness}).\tag{9.1.12.2}
\]
（使用 7.6 h)、e)。）

b) 证明\Term{finite-topos}{有限拓扑斯}是\Term{noetherian-topos}{诺特的}（第 VI 讲 2）。

c) 证明：\Term{finite-topos}{有限拓扑斯} $E$ 的每个\Term{subtopos}{子拓扑斯}都是\Term{finite-topos}{有限拓扑斯}。更确切地，若 $E$ \Term{equivalence}{等价}于 $\widehat C$，其中 $C$ 是\Term{finite}{有限}的\Term{karoubian-category}{Karoubi 范畴}（或更一般地，\Term{equivalence}{等价}于\Term{finite-category}{有限范畴}），证明有如下\Term{order-isomorphism}{有序集合同构}：一方是 $E$ 的\Term{subtopos}{子拓扑斯} $E'$ 的\Term{set}{集合}，另一方是 $C$ 中为\Term{karoubian-category}{Karoubi 范畴}（\Term{equivalence}{等价}地，在 $C$ 中对\Term{direct-factor}{直因子}\Term{stable}{稳定}）的\Term{strictly-full-subcategory}{严格满子范畴} $C'$ 的\Term{set}{集合}；对每个 $C'$，配以 $f':\widehat{C'}\to\widehat C$ 的\Term{essential-image}{本质像}，其中 $f:C'\to C$ 是\Term{inclusion-functor}{包含函子}。（使用 5.6 确认 $f_*$ \Term{fully-faithful-functor}{全忠实}；使用 b) 和 9.1.11 b) 证明 $E$ 的每个\Term{subtopos}{子拓扑斯} $E'$ 都可如上得到。）

d) 设 $E$ 是\Term{finite-topos}{有限拓扑斯}。在 $E$ 的点的\Term{isomorphism}{同构}类所成的\Term{finite}{有限}\Term{set}{集合} $X=\operatorname{Point}(E)$ 上构造一个\Term{topology}{拓扑}，使其成为\Term{sober-space}{清醒}\Term{topological-space}{拓扑空间}，且 $E$ 的\Term{subtopos}{子拓扑斯}所成\Term{ordered-set}{有序集合}\Term{canonical-isomorphism}{典范同构}于 $X$ 的\Term{closed-subset}{闭子集}\Term{set}{集合}。（取 $X$ 的“\Term{constructible-topology}{可构造拓扑}”$\mathscr C_{\mathrm{cons}}$，由\Term{order-relation}{序关系} $(x\le y)\Leftrightarrow(x\text{ 是 }y\text{ 的直因子})$ 定义；其中点 $x$ 的\Term{closure}{闭包}由所有 $y\le x$ 的 $y$ 组成。）特别地，$E$ 的\Term{subtopos}{子拓扑斯}\Term{set}{集合} $\Sigma$ \Term{finite}{有限}，且其中的\Term{infimum}{下确界}对任意\Term{supremum}{上确界}满足分配律。

\SourcePage{272}
e) 设 $C$ 是\Term{equivalence}{等价}于\Term{finite-category}{有限范畴}的\Term{category}{范畴}。证明：对 $C$ 的每个\Term{full-subcategory}{满子范畴} $C'$，$C$ 上存在\Term{topology}{拓扑} $T_{C'}$，使族 $(X_i\to X)_{i\in I}$ 为\Term{covering-family}{覆盖族}，当且仅当对每个 $Y\in\operatorname{Ob}C'$，映射族 $\operatorname{Hom}(Y,X_i)\to\operatorname{Hom}(Y,X)$（$i\in I$）\Term{surjectivity}{满射}。证明：$(C,T_{C'})$ 上的\Term{sheaf}{层}\Term{category}{范畴}\Term{equivalence}{等价}于 $\widehat{C'}$；且 $C'\mapsto T_{C'}$ 在 $C$ 中为\Term{karoubian-category}{Karoubi 范畴}（\Term{equivalence}{等价}地，在 $C$ 中对\Term{direct-factor}{直因子}\Term{stable}{稳定}）的\Term{strictly-full-subcategory}{严格满子范畴} $C'$ 的\Term{set}{集合}与\Term{category}{范畴} $C$ 上的\Term{topology}{拓扑}（第 II 讲 1.1）的\Term{set}{集合}之间建立\Term{bijection}{双射}（反转\Term{order-relation}{序关系}）。（使用 c) 和 9.1.2 c)。）特别地，若 $C$ 是一个\Term{site}{景}，其底层\Term{category}{范畴}\Term{equivalence}{等价}于\Term{finite-category}{有限范畴}，则\Term{topos}{拓扑斯} $C^\sim$ 是\Term{finite-topos}{有限拓扑斯}。

f) 设 $f:C'\to C$ 是 $\mathcal U$-\Term{category}{范畴}之间的\Term{functor}{函子}，其中 $C$ \Term{equivalence}{等价}于\Term{finite-category}{有限范畴}。证明：\Term{topos-morphism}{拓扑斯态射} $\widehat f:\widehat{C'}\to\widehat C$（4.6）是\Term{conservative-morphism}{保守的}，当且仅当 $C$ 的每个\Term{object}{对象}都\Term{isomorphism}{同构}于 $f$ 的像中某个\Term{object}{对象}的\Term{direct-factor}{直因子}。（归结到 $f$ 是包含 $f:C'\hookrightarrow C$ 且 $C'$ 如 c) 所述的情形。）

g) 设 $C$ 是如下\Term{category}{范畴}：它有两个\Term{object}{对象} $e$（\Term{final-object}{终对象}）和 $a$，除恒等\Term{arrow}{箭头}外，还有三个\Term{arrow}{箭头} $f:a\to e$、$g:e\to a$、$p=gf$（故 $p^2=p$）。证明：$E=\widehat C$ 除\Term{subtopos}{子拓扑斯} $E$ 和 $\operatorname{Top}(\varnothing)$ 外，恰好还有一个\Term{subtopos}{子拓扑斯} $E'$，即 $E'=\widehat{\{e\}}$；从而对 $E$ 的每个\Term{subtopos}{子拓扑斯} $E''\ne\operatorname{Top}(\varnothing)$，都有 $E'\cap E''\ne\operatorname{Top}(\varnothing)$。
\end{SGAresult}
\begin{SGAresult}{练习}{9.1.13}
（\Term{topos}{拓扑斯}的\Term{complementary-subtopos}{互补子拓扑斯}。）设 $E$ 是\Term{topos}{拓扑斯}，$E',E''$ 是其两个\Term{subtopos}{子拓扑斯}。若 $E'\cap E''=\operatorname{Top}(\varnothing)$ 且 $E'\vee E''=E$，则称 $E'$ 和 $E''$ \Term{complementary-subtopos}{互补}。这里 $\vee$ 表示 $E$ 的\Term{subtopos}{子拓扑斯}所成\Term{ordered-set}{有序集合}中的\Term{supremum}{上确界}。

假设 $E$ 的\Term{subtopos}{子拓扑斯}\Term{set}{集合} $\Sigma$ 中，二元\Term{infimum}{下确界}对\Term{finite}{有限}\Term{supremum}{上确界}满足分配律。

a) 若 $E'$ 和 $E''$ 互补，则 $E''$ 是满足 $E'\cap E''_1=\operatorname{Top}(\varnothing)$ 的\Term{subtopos}{子拓扑斯} $E''_1$ 中最大的一个（即按 9.1.10 d) 的术语，$E''$ 是 $E'$ 的\Term{weak-complement-subtopos}{弱补}）。

b) 证明 $E$ 的下列条件\Term{equivalence}{等价}：
\begin{enumerate}
\item[(i)] 若 $E',E''$ 是 $E$ 的两个\Term{subtopos}{子拓扑斯}，且 $E''$ 是 $E'$ 的\Term{weak-complement-subtopos}{弱补}，则 $E''$ 是 $E'$ 的\Term{complementary-subtopos}{补}。
\item[(ii)] 对每个\Term{subtopos}{子拓扑斯} $E'\ne E$，存在\Term{subtopos}{子拓扑斯} $E''\ne\operatorname{Top}(\varnothing)$，使得 $E'\cap E''=\operatorname{Top}(\varnothing)$。
\item[(i bis)] 若 $E',E''$ 是 $E$ 的两个\Term{subtopos}{子拓扑斯}，且 $E''$ 在满足 $E'\cap E''_1=\operatorname{Top}(\varnothing)$ 的\Term{subtopos}{子拓扑斯} $E''_1$ 的\Term{set}{集合}中为极大元素，则 $E''$ 是 $E$ 的\Term{complementary-subtopos}{补}。
\end{enumerate}
证明即使 $E$ 是\Term{finite-topos}{有限拓扑斯}（9.1.12），这些条件也未必满足（9.1.12 g)）。不过，若 $E=\operatorname{Top}(X)$，其中 $X$ 是\Term{locally-noetherian-space}{局部诺特空间}，则它们成立。（使用 9.1.11 b)。）

\SourcePage{273}
c) 若 $E$ 的\Term{subtopos}{子拓扑斯} $E'$ 有一个\Term{complementary-subtopos}{补}，则称它为\Term{complemented-subtopos}{有补的}（比较 9.1.12 c)）。证明：有补的\Term{subtopos}{子拓扑斯}\Term{set}{集合}对\Term{finite}{有限}\Term{infimum}{下确界}和\Term{supremum}{上确界}以及取补运算\Term{stable}{稳定}；取补把\Term{infimum}{下确界}变为\Term{supremum}{上确界}，把\Term{supremum}{上确界}变为\Term{infimum}{下确界}。

d) 陈述 a)、b)、c) 的\Term{duality}{对偶}观察（它们事实上是关于抽象\Term{ordered-set}{有序集合} $\Sigma$ 的\Term{trivial-bundle}{平凡}一般命题；注意对 $\Sigma$ 所作的假设本身是\Term{autoduality}{自对偶}的）。
\end{SGAresult}
\textbf{9.1.14. 开放问题。}设 $E$ 是一个\Term{topos}{拓扑斯}。

a) 在 $E$ 的\Term{subtopos}{子拓扑斯}\Term{set}{集合}中，二元\Term{infimum}{下确界}对\Term{finite}{有限}\Term{supremum}{上确界}是否满足分配律？

b) 设 $E',E''$ 是 $E$ 的两个\Term{subtopos}{子拓扑斯}，且 $E=E'\vee E''$。是否 $E$ 的每个点都\Term{isomorphism}{同构}于 $E'$ 或 $E''$ 的某个点的像？

注意：若 b) 对 $E$ 及其所有\Term{subtopos}{子拓扑斯}都有肯定回答，则 a) 至少对有\Term{enough-points}{足够多的点}的\Term{subtopos}{子拓扑斯}有肯定回答（因此，若 $E$ \Term{locally-noetherian}{局部诺特}，则对所有\Term{subtopos}{子拓扑斯}都有肯定回答，9.1.11 b)）。然而，对一切\Term{noetherian-topos}{诺特拓扑斯}，a) 的回答（因而更不用说 b) 的回答）仍未知。不过，当 $E=\operatorname{Top}(X)$，$X$ 是\Term{locally-noetherian-space}{局部诺特空间}时（9.1.11），或当 $E$ 是\Term{finite-topos}{有限拓扑斯}时（9.1.12），a)、b) 的回答都是肯定的。

注意 a) 可以改述如下（对 $E$ 的所有\Term{subtopos}{子拓扑斯} $F$ 应用）：

a') 若 $E',E''$ 是\Term{topos}{拓扑斯} $F$ 的两个\Term{subtopos}{子拓扑斯}，且\Term{canonical}{典范}\Term{morphism}{态射}（参见 8.7 b)）
\[
E'\amalg E''\longrightarrow F\tag{9.1.14.1}
\]
是\Term{conservative-morphism}{保守的}（这也表示 $E=E'\vee E''$，9.1.7.2 e)），则经 $2$-\Term{base-change}{基变换} $F_1\hookrightarrow F$ 得到的\Term{morphism}{态射}也\Term{conservative-morphism}{保守}，其中 $F_1$ 是 $F$ 的\Term{subtopos}{子拓扑斯}。

取\Term{punctual-topos}{点拓扑斯} $F_1$（2.2）给出的\Term{base-change}{基变换} $F_1\to F$，便得到 b) 的类似改述 b')。于是，a)、b) 的肯定回答将由下列问题的肯定回答推出（其中 $F$ 遍取 $E$ 的\Term{subtopos}{子拓扑斯}）：

c) 若 $E',E''$ 是\Term{topos}{拓扑斯} $F$ 的两个\Term{subtopos}{子拓扑斯}，其\Term{supremum}{上确界}为 $F$，即\Term{topos-morphism}{拓扑斯态射} (9.1.14.1) 是\Term{conservative-morphism}{保守的}，那么这个\Term{morphism}{态射}是否甚至\Term{universally-conservative-morphism}{普遍保守}，即在任意 $2$-\Term{base-change}{基变换} $F_1\to F$ 后仍\Term{conservative-morphism}{保守}？

由 Deligne 关于\Term{topos}{拓扑斯}的 $2$-\Term{fiber-product}{纤维积}存在的结果（\,），这个命题有意义。但参见例 9.1.7.2 d)。

另一方面，由 9.1.11 h) 可见：b) 的肯定回答将在\Term{locally-noetherian}{局部诺特}情形下，为下列问题提供肯定回答：
\SourcePage{274}
d) 在 $E$ 的\Term{subtopos}{子拓扑斯}所成的\Term{ordered-set}{有序集合} $\Sigma$ 中，二元\Term{supremum}{上确界}对任意\Term{infimum}{下确界}（不必\Term{finite}{有限}）是否满足分配律？

这也表示，$\Sigma$ \Term{solvable}{可解}释为某个适当\Term{topos}{拓扑斯} $E_{\mathrm{cons}'}$ 的“\Term{closed-subtopos}{闭子拓扑斯}”所成的\Term{ordered-set}{有序集合}（即 $\Sigma^{0\sim}$，其中 $\Sigma$ 的反序\Term{ordered-set}{有序集合} $\Sigma^0$ 被视为一个\Term{category}{范畴}，赋以其\Term{canonical-topology}{典范拓扑}，参见 7.8 b)）；这个\Term{topos}{拓扑斯}由\Term{final-object}{终对象}的\Term{subobject}{子对象}生成，因此十分接近通常的\Term{topological-space}{拓扑空间}。尚须仿照例 9.1.11 g)、9.1.12 d) 研究这个\Term{topos}{拓扑斯} $E_{\mathrm{cons}'}$，尤其要确定：若 $E$ \Term{noetherian}{诺特}，它是否也\Term{noetherian}{诺特}？这\Term{equivalence}{等价}于下列关于 $E$ 的\Term{subtopos}{子拓扑斯} $F$ 的问题；这个问题独立于 d) 也有意义：

e) 设 $F$ 是\Term{noetherian-topos}{诺特拓扑斯}，$(E_i)_{i\in I}$ 是 $E$ 的一族\Term{subtopos}{子拓扑斯}，其交为 $\operatorname{Top}(\varnothing)$。是否存在一个\Term{finite}{有限}子族，具有同样的性质？

另就一个颇为不同的方向而言，也提醒读者注意 (9.1.6 e)) 中提出的问题，它值得澄清。

\SGAsubsection{9.2}{开子拓扑斯的情形}
\textbf{9.2.1.} 设 $U$ 是\Term{topos}{拓扑斯} $E$ 的\Term{open-object}{开对象}，即 $E$ 的\Term{final-object}{终对象} $e$ 的一个\Term{subobject}{子对象}（8.3）。考虑\Term{localization-morphism}{局部化态射}（5.2）
\[
j:E_{/U}\longrightarrow U.\tag{9.2.2}
\]
因这里 $U\to e$ 是\Term{monomorphism}{单态射}，\Term{functor}{函子} $j_!$（\Term{solvable}{可解}释为遗忘\Term{functor}{函子}）是\Term{fully-faithful-functor}{全忠实的}，故它的双\Term{adjunction}{伴随} $j_*$ 也是\Term{fully-faithful-functor}{全忠实的}（1.5.7 a)）。换言之，\Term{topos}{拓扑斯} $E$ 的\Term{open-object}{开对象} $U$ \Term{correspondence}{对应}的\Term{localization-morphism}{局部化态射} $j$（9.2.2）是\Term{topos-embedding}{拓扑斯嵌入}（9.1.1 a)）。

若一个\Term{topos-embedding}{拓扑斯嵌入} $F\to E$ \Term{isomorphism}{同构}于由 $E$ 的某个\Term{open-object}{开对象} $U$ 定义的\Term{topos-embedding}{拓扑斯嵌入}，则称它为\Term{open-embedding}{开嵌入}。若 $E$ 的\Term{subtopos}{子拓扑斯} $E'$（9.1.1 b)）由 $E$ 的一个\Term{open-object}{开对象}定义，即\Term{canonical}{典范}\Term{inclusion-morphism}{包含态射} $E'\to E$ 是\Term{open-embedding}{开嵌入}，则称 $E'$ 为\Term{open-subtopos}{开子拓扑斯}。注意：\Term{correspondence}{对应}于\Term{topos}{拓扑斯}的\Term{open-embedding}{开嵌入} $j:E'\to E$ 的\Term{open-object}{开对象} $U$ 唯一确定，等于 $e$ 的\Term{subobject}{子对象} $j_!(e')$，其中 $e'$ 表示 $E'$ 的\Term{final-object}{终对象}。于是，$E$ 的\Term{open-object}{开对象} $U$ 与 $E$ 的\Term{open-subtopos}{开子拓扑斯}一一\Term{correspondence}{对应}。

\begin{SGAresult}{注记}{9.2.3}
按前述定义，\Term{correspondence}{对应}于 $E$ 的\Term{open-object}{开对象} $U$ 的\Term{open-subtopos}{开子拓扑斯}，是\Term{direct-image-functor}{正像函子}
\[
j_*:E_{/U}\longrightarrow E
\]
的\Term{essential-image}{本质像}所成的 $E$ 的\Term{strictly-full-subcategory}{严格满子范畴}。须注意不要把 $E$ 的这个\Term{subcategory}{子范畴}，与 $j_!$ 的\Term{essential-image}{本质像}所成的\Term{strictly-full-subcategory}{严格满子范畴}混淆；后者由\Term{structural-morphism}{结构态射} $X\to e$ 可经 $U$ \Term{resolution}{分解}的 $E$ 的\Term{object}{对象} $X$ 组成。
\SourcePage{275}
不过，显然这两个\Term{subcategory}{子范畴}彼此确定，且\Term{canonical}{典范}\Term{equivalence}{等价}。因此有些作者可能倾向于（甚至已经）将 $j_!$ 的\Term{essential-image}{本质像}所成的\Term{strictly-full-subcategory}{严格满子范畴}称为“由 $U$ 定义的\Term{open-subtopos}{开子拓扑斯}”，因为用 $U$ 描述它，比描述 $j_*$ 的\Term{essential-image}{本质像}更简单。只要不打算研究\Term{open-subtopos}{开子拓扑斯}以外的其他种类的\Term{subtopos}{子拓扑斯}，这种术语便没有不便。我们的情形并非如此，故这里不采用上述作者的用法。
\end{SGAresult}
为方便读者，下面汇总\Term{open-embedding}{开嵌入}最重要的特殊性质：
\begin{SGAresult}{命题}{9.2.4}
设 $j:\mathcal U\to E$ 是\Term{topos}{拓扑斯}的\Term{open-embedding}{开嵌入}（9.2.1）。则 $j^*$ 的\Term{left-adjoint}{左伴随}\Term{functor}{函子} $j_!$ 存在，故有三个\Term{adjoint-functor}{伴随函子}组成的序列
\[
j_!,\quad j^*,\quad j_*:
\qquad
\begin{tikzcd}[column sep=large]
\mathcal U\arrow[r,shift left=2.2ex,"j^*"]\arrow[r,shift right=2.2ex,"j_*"']&E\arrow[l,"j_!" description]
\end{tikzcd}
\]
此外：

a) \Term{functor}{函子} $j_!$ 与 $j_*$ 都\Term{fully-faithful-functor}{全忠实}。

b) \Term{functor}{函子} $j_!$ 与\Term{fiber-product}{纤维积}、以\Term{nonempty}{非空}\Term{small-set}{小集合} $I\ne\varnothing$ 为指标的积，以及相对于小\Term{cofiltered-category}{余滤过}指标\Term{category}{范畴}的\Term{projective-limit}{投射极限}（1.2.7）交换。

c) 对 $E$ 的每个\Term{object}{对象} $X$，\Term{adjunction-morphism}{伴随态射}
\[
j_!j^*(X)\longrightarrow X
\]
是\Term{monomorphism}{单态射}。
\end{SGAresult}
\begin{proof}
可以假设 $j$ 是由 $E$ 的\Term{open-object}{开对象} $U$ 定义的\Term{localization-morphism}{局部化态射}。于是 a) 只是重述已知事实；b) 来自 $j_!$ 作为遗忘\Term{functor}{函子}的解释以及 $U\to e$ 是\Term{monomorphism}{单态射}这一事实。最后，注意所论\Term{morphism}{态射}正是由包含 $U\to e$ 经\Term{base-change}{基变换} $X\to e$ 得到的\Term{morphism}{态射}，而\Term{base-change}{基变换}把\Term{monomorphism}{单态射}变为\Term{monomorphism}{单态射}，便立即得到 c)。
\end{proof}
\begin{SGAresult}{注记}{9.2.4.1}
可以找到\Term{topos-embedding}{拓扑斯嵌入} $j:E'\to E$，使 $j_!$ 存在，但 $j_!$ 既不与两个\Term{object}{对象}的积交换，也不与\Term{cofiltered-category}{余滤过}\Term{projective-limit}{投射极限}交换（练习 9.5.9 c)）。
\end{SGAresult}
\SourcePage{276}
\textbf{9.2.5.} 现在设 $g:E'\to E$ 是任意\Term{topos-morphism}{拓扑斯态射}，$U'=g^*(U)$ 为 $E$ 的\Term{open-object}{开对象} $U$ 的\Term{inverse-image}{逆像}。在 5.11 已看到，\Term{topos}{拓扑斯}图
\[
\begin{tikzcd}
E'_{/U'}\arrow[r,"j'"]\arrow[d,"g_U"']&E'\arrow[d,"g"]\\
E_{/U}\arrow[r,"j"']&E
\end{tikzcd}
\]
是 $2$-\Term{cartesian-diagram}{笛卡尔的}。特别推出：$E$ 的\Term{open-subtopos}{开子拓扑斯}在 $g$ 下的\Term{inverse-image}{逆像}（参见 9.1.6 d)）是 $E'$ 的\Term{open-subtopos}{开子拓扑斯}；\Term{equivalence}{等价}地，\Term{topos}{拓扑斯}的\Term{open-embedding}{开嵌入}概念（9.2.1）在 $\mathcal U$-\Term{topos}{拓扑斯}（属于给定\Term{universe}{宇宙} $\mathcal V$）的 $2$-\Term{category}{范畴}中的 $2$-\Term{base-change}{基变换}下\Term{stable}{稳定}。

还要注意：\Term{topos-morphism}{拓扑斯态射} $g:E'\to E$ 在\Term{isomorphism}{同构}意义下可经 $j:E_{/U}\to E$ \Term{resolution}{分解}，当且仅当 $j':E'_{/U'}\to E'$ 是\Term{equivalence-of-categories}{范畴等价}；这又表示包含 $U'\hookrightarrow e'$（$e'$ 为 $E'$ 的\Term{final-object}{终对象}）是\Term{isomorphism}{同构}。这样就以简单方式刻画了 9.1.4 所论\Term{functor}{函子}的\Term{essential-image}{本质像}，从而细化了该处的结论。

\textbf{9.2.6.} 将 9.2.5 应用于 $E'=E_{/V}$ 的情形，其中 $V$ 是 $E$ 的另一个\Term{open-object}{开对象}，$g$ 是\Term{localization-morphism}{局部化态射}。则 $U'=U\cap V$，并得到 $E_{/U}$ 与 $E_{/V}$ 在 $E$ 上的 $2$-\Term{fiber-product}{纤维积}可与 $E_{/U\cap V}$ 等同。于是 $E$ 的\Term{finite}{有限}个\Term{open-subtopos}{开子拓扑斯}的交是 $E$ 的\Term{open-subtopos}{开子拓扑斯}；更确切地，映射
\[
U\longmapsto\text{由 }U\text{ 定义的 }E\text{ 的子拓扑斯}\tag{9.2.6.1}
\]
与\Term{finite}{有限}交交换。
\begin{SGAresult}{练习}{9.2.7}
证明映射 (9.2.6.1) 也与任意\Term{supremum}{上确界}交换。（使用 9.1.7.2 e)。）
\end{SGAresult}
\SGAsubsection{9.3}{与开子拓扑斯互补的闭子拓扑斯的构造}
\textbf{9.3.1.} 设 $T$ 是\Term{topological-space}{拓扑空间}，$U$ 是 $T$ 的\Term{open-subset}{开集}，故 $\operatorname{Top}(U)$ 可与 $\operatorname{Top}(T)$ 的一个\Term{open-subtopos}{开子拓扑斯}等同（采用 2.1 的记号）。设 $Y$ 是 $T$ 中与 $U$ 互补的闭\Term{topological-space}{拓扑子空间}。于是可在\Term{equivalence}{等价}意义下将 $\operatorname{Top}(Y)$ 看作 $\operatorname{Top}(T)$ 的一个\Term{subtopos}{子拓扑斯}，即由那些在 $U$ 上的\Term{restriction}{限制}为 $U$ 的\Term{final-topos}{终拓扑斯}的\Term{object}{对象} $F\in\operatorname{Top}(T)$ 组成的\Term{subtopos}{子拓扑斯}。只要给定一个\Term{topos}{拓扑斯}及其\Term{open-object}{开对象} $U$，对 $E=\operatorname{Top}(T)$ 的一个\Term{subcategory}{子范畴}的上述描述就仍有意义；我们将看到，它总是给出 $E$ 的一个\Term{subtopos}{子拓扑斯}。
\SourcePage{277}
此外，在最初考虑的特殊情形中，立即可见（采用练习 9.1.13 引入的术语）$\operatorname{Top}(Y)$ 与 $\operatorname{Top}(U)$ 是互为\Term{complementary-subtopos}{补}的\Term{subtopos}{子拓扑斯}（使用 9.2.5 和 9.1.7.2 e)）。一般情形仍然如此，并将看到，这一性质唯一刻画了所论\Term{subtopos}{子拓扑斯}。因此，从各个方面说，它都应称为与所论\Term{open-object}{开对象} $U$、或与 $U$ 定义的\Term{open-subtopos}{开子拓扑斯}互补的\Term{closed-subtopos}{闭子拓扑斯}。本节将详细构造这个\Term{topos}{拓扑斯} $\mathcal F$ 及\Term{inverse-image-functor}{逆像函子} $i^*:E\to\mathcal F$，而第 9.4 节将阐述其初步性质。

\textbf{9.3.2.} 于是，如 9.2.1，设 $E$ 是\Term{topos}{拓扑斯}，$U$ 是 $E$ 的\Term{open-object}{开对象}，考虑\Term{localization-morphism}{局部化态射}
\[
j:\mathcal U=E_{/U}\longrightarrow E,\tag{9.3.2.1}
\]
它是一个\Term{open-embedding}{开嵌入}。

对 $E$ 的每个\Term{object}{对象} $X$，置
\[
X_{\mathscr C U}=U\amalg_{U\times X}X,\tag{9.3.2.2}
\]
其中\Term{pushout}{纤维和}相对于\Term{canonical-projection}{典范投影}
\[
\operatorname{pr}_1:U\times X\longrightarrow U,\qquad
\operatorname{pr}_2:U\times X\longrightarrow X
\]
形成。因此，$E$ 中有如下依\Term{functor}{函子}方式依赖于 $X$ 的\Term{cocartesian-diagram}{余笛卡尔图}（第 I 讲 10）：
\[
\begin{tikzcd}
U\times X\arrow[r,"\operatorname{pr}_2"]\arrow[d,"\operatorname{pr}_1"']&X\arrow[d]\\
U\arrow[r]&X_{\mathscr C U}
\end{tikzcd}\tag{9.3.2.3}
\]
注意，在 9.3.1 的例子中，$X_{\mathscr C U}$ \Term{canonical-isomorphism}{典范同构}于 $i_*i^*(X)$（其中 $i:Y\to T$ 为包含），且 $X\to X_{\mathscr C U}$ 等同于\Term{adjunction-morphism}{伴随态射}。
\begin{SGAresult}{命题}{9.3.3}
采用前述记号，有如下结论：

a) 对 $E$ 的每个\Term{object}{对象} $X$，下列条件\Term{equivalence}{等价}：
\begin{enumerate}[label=(\roman*)]
\item 存在 $E$ 的\Term{object}{对象} $Y$ 及\Term{isomorphism}{同构} $X\simeq Y_{\mathscr C U}$。
\item \Term{canonical}{典范}\Term{morphism}{态射} $X\xrightarrow{m}X_{\mathscr C U}$ 是\Term{isomorphism}{同构}。
\item \Term{canonical}{典范}\Term{morphism}{态射} $\operatorname{pr}_2:X\times U\to U$ 是\Term{isomorphism}{同构}（即 $j^*(X)$ 是 $E_{/U}$ 的\Term{final-object}{终对象}）。
\end{enumerate}
b) 对 $E$ 中每个\Term{morphism}{态射} $m:E\to X'$，
\SourcePage{278}
下列条件\Term{equivalence}{等价}：
\begin{enumerate}[label=(\roman*')]
\item 图
\[
\begin{tikzcd}
X\arrow[r,"m"]\arrow[d]&X'\arrow[d]\\
X_{\mathscr C U}\arrow[r]&X'_{\mathscr C U}
\end{tikzcd}
\]
是\Term{cartesian-diagram}{笛卡尔的}。
\item \Term{morphism}{态射} $m\times\operatorname{id}_U:X\times U\to X'\times U$ 是\Term{isomorphism}{同构}，即 $j^*(m)$ 是\Term{isomorphism}{同构}。
\end{enumerate}
c) 从 $E$ 到 $E$ 的\Term{functor}{函子} $X\longmapsto X_{\mathscr C U}$ 是\Term{left-exact-functor}{左正合的}，把\Term{epimorphic-family}{满态射族}变为\Term{epimorphic-family}{满态射族}，且与\Term{filtered}{滤过}\Term{inductive-limit}{归纳极限}及\Term{pushout}{纤维和}交换。
\end{SGAresult}
\begin{proof}
先在 $E=\widehat C$、$C$ 是小\Term{category}{范畴}时证明；为此，由第 I 讲 3.1 可归结到 $E$ 是“\Term{punctual-topos}{点}”\Term{topos}{拓扑斯} $(\mathrm{Ens})$ 的情形，此时命题显然。再借第 II 讲 4 的标准方法，使用“\Term{associated-sheaf}{关联层}”\Term{functor}{函子}，由此过渡到一般情形。
\end{proof}
\begin{SGAresult}{命题}{9.3.4}
采用 9.3.2 的记号，由满足 $j^*(X)$ 为 $\mathcal U$ 的\Term{final-object}{终对象}的 $E$ 的\Term{subobject}{子对象} $X$（参见 9.3.3 a)）定义的 $E$ 的\Term{strictly-full-subcategory}{严格满子范畴} $\mathcal F$，是 $E$ 的一个\Term{subtopos}{子拓扑斯}，即依 9.1.1，它是\Term{topos}{拓扑斯}，且\Term{inclusion-functor}{包含函子} $i_*:\mathcal F\to E$ 是某个\Term{topos-morphism}{拓扑斯态射} $i:\mathcal F\to E$ 的\Term{direct-image-functor}{正像函子}，其\Term{inverse-image-functor}{逆像函子}为 $X\mapsto i^*X=X_{\mathscr C U}$。\Term{adjunction-morphism}{伴随态射} $X\to i_*i^*X$ 是 (9.3.2.3) 中的\Term{canonical}{典范}\Term{morphism}{态射} $X\to X_{\mathscr C U}$。
\end{SGAresult}
\begin{proof}
对 $E$ 的每个\Term{object}{对象} $X$，以 $u(X)$ 表示\Term{canonical}{典范}\Term{morphism}{态射} $X\to X_{\mathscr C U}$。\Term{morphism}{态射} $u(X)$ 关于 $X$ 是\Term{functorial}{函子性}的，且对每个\Term{object}{对象} $X$，$u(X_{\mathscr C U})$ 是\Term{isomorphism}{同构}。由这两条性质，形式上立即得到：\Term{functor}{函子} $i^*$ 是\Term{functor}{函子} $i_*$ 的\Term{left-adjoint}{左伴随}，且\Term{adjunction-morphism}{伴随态射}为 $u(X)$。因\Term{functor}{函子} $X\mapsto X_{\mathscr C U}$ 是\Term{left-exact-functor}{左正合的}，\Term{functor}{函子} $i^*$ 也是\Term{left-exact-functor}{左正合的}。由第 II 讲 5.5，$\mathcal F$ 是\Term{topos}{拓扑斯}；由定义 3.1，$(i^*,i_*)$ 是\Term{topos-morphism}{拓扑斯态射}。
\end{proof}
\textbf{9.3.5.} 9.3.4 描述的 $E$ 的\Term{subtopos}{子拓扑斯} $\mathcal F$，称为 $E$ 中与\Term{open-object}{开对象} $U$ 互补的\Term{closed-subtopos}{闭子拓扑斯}，或也可说，与\Term{correspondence}{对应}于 $U$ 的\Term{open-subtopos}{开子拓扑斯}互补的\Term{closed-subtopos}{闭子拓扑斯}。按 9.1.1 b)，9.3.4 构造的\Term{morphism}{态射} $i:\mathcal F\to E$ 称为\Term{inclusion-morphism}{包含态射}。若存在 $E$ 的\Term{open-object}{开对象} $U$，使 $\mathcal F$ 为与 $U$ 互补的\Term{closed-subtopos}{闭子拓扑斯}，则称 $E$ 的\Term{full-subcategory}{满子范畴} $\mathcal F$ 为 $E$ 的\Term{closed-subtopos}{闭子拓扑斯}。
\SourcePage{279}
注意：这个 $U$ 唯一确定，为 $i_*(\varnothing_{\mathcal F})$，其中 $\varnothing_{\mathcal F}$ 是 $\mathcal F$ 的\Term{initial-object}{始对象}，等于 $i^*(\varnothing_E)$；也称 $U$ 为 $E$ 中与\Term{closed-subtopos}{闭子拓扑斯} $\mathcal F$ 互补的\Term{open-object}{开对象}，而 $U$ 定义的 $E$ 的\Term{subtopos}{子拓扑斯}也称为 $E$ 中与 $\mathcal F$ 互补的\Term{open-subtopos}{开子拓扑斯}。最后，若\Term{topos-embedding}{拓扑斯嵌入}（9.1.1 a)）$i:F\to E$ \Term{correspondence}{对应}的 $E$ 的\Term{subtopos}{子拓扑斯}是闭的，则称 $i$ 为\Term{closed-embedding}{闭嵌入}。

若 $G$ 是 $E$ 上的\Term{group-object}{群}，$s\in\operatorname{Hom}(e_E,G)$，则称与 $G$（相应地，$s$）的开\Term{cosupport}{余支撑}（8.5.1、8.5.2）互补的 $E$ 的\Term{closed-subtopos}{闭子拓扑斯}为 $G$（相应地，$S$）的\Term{support}{支撑}。

\textbf{9.3.6.} 采用 9.3.2、9.3.4 的记号，得到两个\Term{topos-morphism}{拓扑斯态射}
\[
\mathcal U\xrightarrow{j}E\xleftarrow{i}\mathcal F,\tag{9.3.6.1}
\]
定义了五个\Term{functor}{函子}，组成两个\Term{adjoint-functor}{伴随函子}序列 $(j_!,j^*,j_*)$ 和 $(i^*,i_*)$：
\[
\begin{tikzcd}[column sep=large]
\mathcal U\arrow[r,shift left=2.2ex,"j_!"]\arrow[r,shift right=2.2ex,"j_*"']&E\arrow[l,"j^*" description]\arrow[r,shift left=1.2ex,"i^*"]&\mathcal F\arrow[l,shift left=1.2ex,"i_*"]
\end{tikzcd}\tag{9.3.6.2}
\]
\Term{functor}{函子}
\[
\rho=i_*j^*:\mathcal U\longrightarrow\mathcal F\tag{9.3.6.3}
\]
称为关于 $E$ 的\Term{open-object}{开对象} $U$ 的\Term{gluing-functor}{粘合函子}；更一般地，若 $\mathcal U,\mathcal F$ 是任意两个\Term{topos}{拓扑斯}，则称\Term{functor}{函子} $\rho:\mathcal U\to\mathcal F$ 为\Term{gluing-functor}{粘合函子}，若它\Term{left-exact-functor}{左正合}且\Term{accessible-functor}{可达}（第 I 讲 9.2），或\Term{equivalence}{等价}地（1.8），它有一个\Term{pro-adjoint}{pro-伴随}。这一术语的理由将在下文 9.4.1 d) 和 9.5.4 中显现。

\SGAsubsection{9.4}{闭子拓扑斯 $\mathcal F$ 及 $i:\mathcal F\to E$ 的初步性质}
\begin{SGAresult}{命题}{9.4.1}
采用 9.3.2、9.3.4 的记号。

a) \Term{functor}{函子} $i_*:\mathcal F\to E$ 把\Term{epimorphic-family}{满态射族}变为\Term{epimorphic-family}{满态射族}。它与\Term{pushout}{纤维和}及\Term{filtered}{滤过}\Term{inductive-limit}{归纳极限}的形成交换。

b) 对 $E$ 的每个满足\Term{projection}{投影}\Term{morphism}{态射} $U\times X\to U$ 为\Term{epimorphism}{满态射}的\Term{object}{对象} $X$，\Term{adjunction-morphism}{伴随态射} $X\to i_*i^*X$ 是\Term{epimorphism}{满态射}。

c) 分别从 $E$ 到 $\mathcal F$、$\mathcal U$ 的\Term{functor}{函子}对 $(i^*,j^*)$ 是\Term{conservative-family}{保守的}（第 I 讲 6.1），即 $E$ 中\Term{morphism}{态射} $m$ 是\Term{isomorphism}{同构}，当且仅当 $i^*(m)$ 和 $j^*(m)$ 都是\Term{isomorphism}{同构}（这\Term{equivalence}{等价}于\Term{functor}{函子}对 $(i^*,j^*)$ 是\Term{faithful-functor}{忠实的}，6.4.0）。

d) \Term{gluing-functor}{粘合函子} $\rho=i^*j_*:\mathcal U\to\mathcal F$ 是\Term{left-exact-functor}{左正合}且\Term{accessible-functor}{可达的函子}（第 I 讲 9.2），或\Term{equivalence}{等价}地（第 I 讲 8），它有一个\Term{pro-adjoint}{pro-伴随} $\sigma:\operatorname{Pro}\mathcal F\to\operatorname{Pro}\mathcal U$。
\end{SGAresult}
\SourcePage{280}
\begin{proof}
a) 由 9.3.3 c) 得到。证明 b) 时，只须参看 $X_{\mathscr C U}$ 的定义 (9.3.2.2)，并考虑 $i_*i^*X=X_{\mathscr C U}$（9.3.4）。断言 c) 由 9.3.3 b) 得到。\Term{functor}{函子} $i^*$、$j_*$ 都\Term{left-exact-functor}{左正合}，故 $\rho$ \Term{left-exact-functor}{左正合}。\Term{functor}{函子} $j_*$ 是某个\Term{topos-morphism}{拓扑斯态射}的\Term{direct-image-functor}{正像函子}，故\Term{accessible-functor}{可达}（第 I 讲 9.5）；注意\Term{topos}{拓扑斯}是\Term{accessible-category}{可达范畴}（第 I 讲 9.11.3）。又因\Term{functor}{函子} $i^*$ 与\Term{inductive-limit}{归纳极限}交换（它是\Term{inverse-image-functor}{逆像函子}），$\rho$ 作为两个\Term{accessible-functor}{可达函子}的复合，亦\Term{accessible-functor}{可达}。
\end{proof}
\begin{SGAresult}{命题}{9.4.2}
设 $E'$ 是\Term{topos}{拓扑斯}。则\Term{functor}{函子} $h\mapsto i\circ h$：
\[
\operatorname{Homtop}(E',\mathcal F)\longrightarrow\operatorname{Homtop}(E',E)
\]
是\Term{fully-faithful-functor}{全忠实的}，其\Term{essential-image}{本质像}由满足 $g^*(U)$ 是 $E'$ 的\Term{initial-object}{始对象} $\varnothing_{E'}$ 的\Term{topos-morphism}{拓扑斯态射} $g:E'\to E$ 组成。
\end{SGAresult}
\begin{proof}
由 9.1.4，知道所论\Term{functor}{函子}是\Term{fully-faithful-functor}{全忠实的}，且 $g$ 属于其\Term{essential-image}{本质像}，当且仅当对 $E'$ 的每个\Term{object}{对象} $X'$，$E$ 的\Term{object}{对象} $g_*(X')$ 都属于 $\mathcal F$，即
\[
j^*(g_*(X'))\simeq e_{\mathcal U}\quad\text{对每个 }X'\in\operatorname{Ob}E'.\tag{*}
\]
置 $U'=g^*(U)$，并考虑\Term{topos}{拓扑斯}图
\[
\begin{tikzcd}
E'_{/U'}\arrow[r,"j'"]\arrow[d,"g_U"']&E'\arrow[d,"g"]\\
E_{/U}\arrow[r,"i"']&E;
\end{tikzcd}
\]
显然有\Term{canonical}{典范}\Term{isomorphism}{同构}
\[
j^*(g_*(X'))\simeq g_{U*}(j'^*(X')).
\]
另一方面，\Term{functor}{函子} $j'^*$ 显然\Term{essentially-surjective-functor}{本质满射}，所以条件 (*) \Term{equivalence}{等价}于
\[
g_{U*}(Y')\simeq e_{\mathcal U}\quad\text{对每个 }Y'\in\operatorname{Ob}E'_{/U'}.\tag{**}
\]
使用 $g_{U*}$ 与 $g_U^*$ 的\Term{adjunction}{伴随}性质，立刻看出：这表示对 $E_{/U}$ 的每个\Term{object}{对象} $Y$，$g_U^*(Y)$ 都是 $E'_{/U'}$ 的\Term{initial-object}{始对象}；或者，利用 $E'_{/U'}$ 的\Term{initial-object}{始对象}是严格的（第 II 讲 4.5），\Term{equivalence}{等价}于取 $Y=e_{\mathcal U}$，即 $\mathcal U$ 的\Term{final-object}{终对象}时有此性质。但这也表示 $E'_{/U'}$ 的\Term{final-object}{终对象} $U'$ 是\Term{initial-object}{始对象}；而这又显然\Term{equivalence}{等价}于 $U'$ 是 $E'$ 的\Term{initial-object}{始对象} $\varnothing_{E'}$。
\end{proof}
\begin{SGAresult}{推论}{9.4.3}
设 $g:E'\to E$ 是\Term{topos-morphism}{拓扑斯态射}，$U'=g^*(U)$，$\mathcal U'=E'_{/U'}$，$\mathcal F'$ 是 $E'$ 中与\Term{open-object}{开对象} $U'$ 互补的\Term{closed-subtopos}{闭子拓扑斯}，$i':\mathcal F'\to E'$ 为\Term{inclusion-morphism}{包含态射}。则\Term{topos-morphism}{拓扑斯态射} $g\cdot i':\mathcal F\to E$ 在\Term{isomorphism}{同构}意义下经由\Term{topos-morphism}{拓扑斯态射} $g_{\mathcal F}:\mathcal F'\to\mathcal F$ \Term{resolution}{分解}，
\SourcePage{281}
且相应的\Term{topos}{拓扑斯}图
\[
\begin{tikzcd}
\mathcal F'\arrow[r,"i'"]\arrow[d,"g_{\mathcal F}"']\arrow[dr,Rightarrow,shorten <=18pt,shorten >=18pt]&E'\arrow[d,"g"]\\
\mathcal F\arrow[r,"i"']&E
\end{tikzcd}\tag{9.4.3.1}
\]
是 $2$-\Term{cartesian-diagram}{笛卡尔的}（参见 5.11）。
\end{SGAresult}
这由 9.4.2 及定义形式地得到。

采用 9.1.6 d) 引入的术语，特别可以说：$E$ 的\Term{closed-subtopos}{闭子拓扑斯} $\mathcal F$ 在\Term{topos-morphism}{拓扑斯态射} $g:E'\to E$ 下的\Term{inverse-image}{逆像}，是 $E'$ 的\Term{closed-subtopos}{闭子拓扑斯} $\mathcal F'$（即与\Term{open-object}{开对象} $U'=g^*(U)$ 互补的\Term{closed-subtopos}{闭子拓扑斯}，其中 $U$ 是 $E$ 中与 $\mathcal F$ 互补的\Term{open-object}{开对象}）。
\begin{SGAresult}{推论}{9.4.4}
设 $E$ 是\Term{topos}{拓扑斯}，$\mathcal U$ 是 $E$ 的\Term{open-subtopos}{开子拓扑斯}，$\mathcal F$ 是互补的\Term{closed-subtopos}{闭子拓扑斯}。

a) 有
\[
\mathcal U\cap\mathcal F\simeq\operatorname{Top}(\varnothing),\qquad
\mathcal U\vee\mathcal F=E,
\]
其中 $\vee$ 表示 $E$ 的所有\Term{subtopos}{子拓扑斯}所成\Term{set}{集合}中的\Term{supremum}{上确界}（9.1.2 c)）。（按 9.1.13 引入的术语，$\mathcal U$ 与 $\mathcal F$ 是 $E$ 中互补的\Term{subtopos}{子拓扑斯}。）

b) $\mathcal U$（相应地，$\mathcal F$）是满足 $E'\cap\mathcal F\simeq\operatorname{Top}(\varnothing)$（相应地，$\mathcal U\cap E'\simeq\operatorname{Top}(\varnothing)$）的\Term{subtopos}{子拓扑斯} $E'$ 中最大的一个。
\end{SGAresult}
\begin{proof}
a) 第一个关系\Term{equivalence}{等价}于说：若 $E'$ 是\Term{topos}{拓扑斯}，且\Term{topos-morphism}{拓扑斯态射} $g:E'\to E$ 在\Term{isomorphism}{同构}意义下既可经 $\mathcal U$ 又可经 $\mathcal F$ \Term{resolution}{分解}，则 $E'\simeq\operatorname{Top}(\varnothing)$。由 9.2.5、9.4.2，假设表示 $g^*(U)$ 同时是 $E'$ 的\Term{initial-object}{始对象}与\Term{final-object}{终对象}，故得结论。对第二个关系，注意由 9.1.7.2 e)，它\Term{equivalence}{等价}于 9.4.1 c)。

b) 假设 $E'\cap\mathcal F\simeq\operatorname{Top}(\varnothing)$；由 9.4.3，$E'\cap\mathcal F$ \Term{equivalence}{等价}于 $E'$ 中与 $g^*(U)$ 互补的\Term{subtopos}{子拓扑斯}，其中 $g:E'\to E$ 是\Term{inclusion-morphism}{包含态射}；故所论条件表示 $U'=g^*(U)$ 是 $E'$ 的\Term{final-object}{终对象}，即依 9.2.5，$E'\subset\mathcal U$。再假设 $\mathcal U\cap E'\simeq\operatorname{Top}(\varnothing)$；由 5.11，$\mathcal U\cap E'$ \Term{equivalence}{等价}于 $E'_{/U'}$，故所论条件表示 $U'$ 是 $E'$ 的\Term{initial-object}{始对象}，即依 9.4.2，$E'\subset\mathcal F$。
\end{proof}
\begin{SGAresult}{推论}{9.4.5}
设 $(\mathcal F_i)_{i\in I}$ 是\Term{topos}{拓扑斯} $E$ 的一族\Term{closed-subtopos}{闭子拓扑斯}，对每个 $i\in I$，设 $U_i$ 是 $E$ 中与 $\mathcal F_i$ 互补的\Term{open-object}{开对象}。则各 $\mathcal F_i$ 的交\Term{subtopos}{子拓扑斯}（9.1.3）是\Term{closed-subtopos}{闭子拓扑斯}，其互补\Term{open-object}{开对象}为 $U=\sup_iU_i$。
\end{SGAresult}
\SourcePage{282}
这形式地来自交作为 $2$-积的\Term{universal-property}{泛性质}及 9.4.2，并考虑到：对\Term{topos-morphism}{拓扑斯态射} $g:E'\to E$，$g^*(U)=\sup_i g^*(U_i)$ 是\Term{initial-object}{始对象}，当且仅当各 $g^*(U_i)$ 都是\Term{initial-object}{始对象}。

类似地证明：
\begin{SGAresult}{推论}{9.4.6}
采用 9.4.5 的记号，假设 $I$ \Term{finite}{有限}。则作为 (9.1.2 c)) 中\Term{supremum}{上确界}的 $E$ 的\Term{subtopos}{子拓扑斯}，是 $E$ 的\Term{closed-subtopos}{闭子拓扑斯}，其互补\Term{open-object}{开对象}等于 $\bigcup_iU_i$。
\end{SGAresult}
\begin{SGAresult}{练习}{9.4.7}
（\Term{subtopos}{子拓扑斯}的\Term{gluing}{粘合}。）设 $E$ 是\Term{topos}{拓扑斯}。

a) 设 $E'$ 是 $E$ 的\Term{subtopos}{子拓扑斯}，$(S_i)_{i\in I}$ 是 $E$ 的一族\Term{object}{对象}，\Term{covering}{覆盖}\Term{final-object}{终对象}；对每个 $i\in I$，设 $S'_i$ 是 $S_i$ 在 $E'$ 中的\Term{inverse-image}{逆像}，$E'_i\simeq E'_{/S'_i}$ 是 $E_i=E_{/S_i}$ 的\Term{subtopos}{子拓扑斯}，为 $E$ 的\Term{subtopos}{子拓扑斯} $E'$ 在\Term{localization-morphism}{局部化态射} $E_i\to E$ 下的\Term{inverse-image}{逆像}（9.1.6 d) 和 5.11）。证明：$E$ 的\Term{object}{对象} $X$ 属于 $E'$，当且仅当对每个 $i\in I$，$X$ 在 $E_i$ 中的\Term{inverse-image}{逆像} $X|S_i$ 属于 $E'_i$。

b) 假设 $E=C^\sim$，其中 $C$ 是 $\mathcal U$-\Term{site}{景}。证明 $C$ 上的\Term{presheaf}{预层}
\[
S\longmapsto\text{拓扑斯 }(C_{/S})^\sim\simeq E^\sim_{/\Sigma(S)}\text{ 的子拓扑斯集合}
\]
是\Term{sheaf}{层}。重新取任意 $E$，由此推出 $E$ 中存在一个元素，表示\Term{functor}{函子}
\[
S\longmapsto\text{拓扑斯 }E_{/S}\text{ 的子拓扑斯集合}.
\]
c) 采用 a) 的记号，证明：$E$ 的\Term{subtopos}{子拓扑斯} $E'$ 是\Term{open-subtopos}{开子拓扑斯}（相应地，\Term{closed-subtopos}{闭子拓扑斯}），当且仅当对每个 $i\in I$，$E_i$ 的诱导\Term{subtopos}{子拓扑斯} $E'_i$ 都如此。
\end{SGAresult}
\begin{SGAresult}{练习}{9.4.8}
（\Term{interior-subtopos}{子拓扑斯的内部}、外部、\Term{closure}{闭包}与\Term{boundary}{边界}。）设 $E$ 是\Term{topos}{拓扑斯}，$E'$ 是\Term{subtopos}{子拓扑斯}，$i:E'\to E$ 是\Term{inclusion-morphism}{包含态射}。

a) 证明 $i_*(\varnothing_{E'})$ 是满足 $E'|\mathcal U\simeq\operatorname{Top}(\varnothing)$ 的最大\Term{open-object}{开对象} $U$，其中 $\mathcal U$ 是由 $U$ 定义的 $E$ 的\Term{open-subtopos}{开子拓扑斯}。称这个 $U$ 或 $\mathcal U$ 为 $E$ 的\Term{subtopos}{子拓扑斯} $E'$ 的\Term{exterior-subtopos}{外部}，称 $E$ 中与 $U$ 互补的\Term{closed-subtopos}{闭子拓扑斯}为 $E'$ 的\Term{closure-subtopos}{闭包}。

b) 证明：存在 $E$ 的最大\Term{open-object}{开对象} $V$，使 $E$ 的相应\Term{open-subtopos}{开子拓扑斯} $\mathcal V$ 包含于 $E'$。（使用 9.2.7。）称这个 $V$ 或 $\mathcal V$ 为 $E$ 的\Term{subtopos}{子拓扑斯} $E'$ 的\Term{interior-subtopos}{内部}。$E$ 的\Term{open-object}{开对象} $U,V$ 不交（$U\cap V=\varnothing_E$）；$E$ 中与 $\sup(U,V)=\mathcal U\amalg\mathcal V$ 互补的\Term{closed-subtopos}{闭子拓扑斯}称为 $E$ 的\Term{subtopos}{子拓扑斯} $E'$ 的\Term{boundary-subtopos}{边界}。

c) $E'$ 是 $E$ 的\Term{open-subtopos}{开子拓扑斯}（相应地，\Term{closed-subtopos}{闭子拓扑斯}），当且仅当它等于其\Term{interior-subtopos}{内部}（相应地，等于其\Term{closure-subtopos}{闭包}）。$E'$ 的\Term{boundary-subtopos}{边界}为空，当且仅当 $E'$ 是 $E$ 中既开又闭的\Term{subtopos}{子拓扑斯}；这又\Term{equivalence}{等价}于它是由 $E$ 的一个\Term{open-object}{开对象} $V$ 定义的\Term{open-subtopos}{开子拓扑斯}，且 $V$ 是\Term{final-object}{终对象} $e$ 的\Term{direct-summand}{直和项}，即存在 $e$ 的\Term{subobject}{子对象} $V$，使 $e\simeq U\amalg V$。
\SourcePage{283}
$E'$ 的\Term{exterior-subtopos}{外部}“为空”，即其\Term{closure-subtopos}{闭包}等于 $E$，当且仅当 $i:E'\to E$ 是\Term{dominant-morphism}{支配态射}（8.8）。

d) 设 $S$ 是 $E$ 的\Term{object}{对象}，$F=E_{/S}$ 为\Term{induced-topos}{诱导拓扑斯}，$F'\simeq F_{/S'}$ 为 $E$ 的\Term{subtopos}{子拓扑斯} $E'$ 所诱导的 $F$ 的\Term{subtopos}{子拓扑斯}。证明：$F'$ 的\Term{exterior-subtopos}{外部}（相应地，\Term{closure-subtopos}{闭包}、\Term{interior-subtopos}{内部}、\Term{boundary-subtopos}{边界}）由 $E'$ 的\Term{exterior-subtopos}{外部}（相应地，……）诱导。
\end{SGAresult}
\begin{SGAresult}{练习}{9.4.9}
（\Term{topos}{拓扑斯}的\Term{locally-closed-subtopos}{局部闭子拓扑斯}。）若\Term{topos}{拓扑斯} $E$ 的\Term{subtopos}{子拓扑斯} $F$ 是一个\Term{open-subtopos}{开子拓扑斯}与一个\Term{closed-subtopos}{闭子拓扑斯}的交，则称 $F$ 为 $E$ 的\Term{locally-closed-subtopos}{局部闭子拓扑斯}。

a) 证明：$E$ 的\Term{finite}{有限}个\Term{locally-closed-subtopos}{局部闭子拓扑斯}的交是\Term{locally-closed-subtopos}{局部闭的}。若 $g:E'\to E$ 是\Term{topos-morphism}{拓扑斯态射}，证明：$E$ 的\Term{locally-closed-subtopos}{局部闭子拓扑斯}的\Term{inverse-image}{逆像}（9.1.6 d)）是 $E'$ 的\Term{locally-closed-subtopos}{局部闭子拓扑斯}。

b) 设 $X$ 是\Term{topological-space}{拓扑空间}，$E=\operatorname{Top}(X)$。对 $X$ 的每个\Term{subset}{子集} $X'$，考虑 $E$ 的\Term{subtopos}{子拓扑斯} $T(X')$，它是 $\operatorname{Top}(X')$ 在 $\operatorname{Top}(X)$ 中的\Term{essential-image}{本质像}。证明：若 $X'$ 是\Term{locally-closed-subset}{局部闭子集}，则 $T(X')$ 是 $E$ 的\Term{locally-closed-subtopos}{局部闭子拓扑斯}；若进一步假设 $X$ 与 $X'$ 清醒，则逆命题成立。证明：$X'\mapsto T(X')$ 在 $X$ 的\Term{locally-closed-subset}{局部闭子集} $X'$ 的\Term{set}{集合}与 $E$ 的\Term{locally-closed-subtopos}{局部闭子拓扑斯}\Term{set}{集合}之间建立\Term{order-isomorphism}{有序集合同构}。

c) 设 $F$ 是 $E$ 的\Term{subtopos}{子拓扑斯}。证明：$F$ 的每个\Term{open-subtopos}{开子拓扑斯}都由 $E$ 的一个\Term{open-subtopos}{开子拓扑斯}诱导。（使用 9.1.8 c)。）证明：$F$ 是 $E$ 的\Term{locally-closed-subtopos}{局部闭子拓扑斯}，当且仅当它是其\Term{closure-subtopos}{闭包}的\Term{open-subtopos}{开子拓扑斯}（9.4.8 a)）。使用 9.4.8 d) 由此推出：$F$ 是 $E$ 的\Term{locally-closed-subtopos}{局部闭子拓扑斯}这一性质，在 $E$ 上是\Term{local-nature}{局部性质}。

d) 设 $F$ 是 $E$ 的\Term{locally-closed-subtopos}{局部闭子拓扑斯}。证明：在将 $F$ 写成 $E$ 的一个\Term{open-subtopos}{开子拓扑斯} $\mathcal U$ 与一个\Term{closed-subtopos}{闭子拓扑斯} $\mathcal F$ 的交的所有方式中，有一种使 $\mathcal U$ 尽可能大，$\mathcal F$ 尽可能小。（取 $\mathcal F$ 为 $F$ 的\Term{closure-subtopos}{闭包}，$\mathcal U$ 为在 $\mathcal F$ 上诱\Term{derivation}{导出}\Term{subtopos}{子拓扑斯} $F$ 的最大\Term{open-subtopos}{开子拓扑斯}。）$\mathcal U,\mathcal F$ 的形成与 $E$ 上的\Term{localization}{局部化}相容。
\end{SGAresult}
\begin{SGAresult}{练习}{9.4.10}
仿照\Term{topological-space}{拓扑空间}情形中熟悉的概念（EGA $0_{\mathrm{III}}$ 9.1，EGA IV 10.1），发展\Term{constructible-subtopos}{可构造子拓扑斯}、\Term{locally-constructible-subtopos}{局部可构造子拓扑斯}、\Term{quasi-constructible-subtopos}{拟可构造子拓扑斯}、\Term{locally-quasi-constructible-subtopos}{局部拟可构造子拓扑斯}的概念。对形如 $\operatorname{Top}(X)$ 的\Term{topos}{拓扑斯}，将重新得到 $X$ 的\Term{constructible-subset}{可构造}（相应地，……）\Term{subset}{子集}的\Term{set}{集合}与 $\operatorname{Top}(X)$ 的\Term{constructible-subtopos}{可构造}（相应地，……）\Term{subtopos}{子拓扑斯}\Term{set}{集合}之间的\Term{bijection}{双射}。
\end{SGAresult}
